\documentclass[12pt]{amsart}
\usepackage[utf8]{inputenc}

\usepackage{comment}
\usepackage[english]{babel}
\usepackage[T1]{fontenc}
\usepackage{todonotes}
\usepackage{bm}
\usepackage{upgreek}


\textheight 22cm
\textwidth 16cm
\topmargin -0.8cm
\oddsidemargin -0.6cm
\evensidemargin -0.6cm



\usepackage{enumitem}
\usepackage{amssymb,amsfonts,stmaryrd,mathabx}
\usepackage{amsmath,latexsym,amsthm,mathtools,bbm}
\usepackage{pgf,tikz} 
\usepackage{subcaption}
\usepackage[colorlinks,cite color=green,link color=purple]{hyperref}
\usepackage{todonotes}
\usepackage{cases}
\usepackage[thinlines]{easytable}
\usepackage{bm}
\usepackage[export]{adjustbox}
\usepackage{appendix}


\allowdisplaybreaks




\usepackage{thmtools}
\usepackage[capitalize,noabbrev,nameinlink]{cleveref}
\newtheorem{thm}{Theorem}[section]
\newtheorem{cor}[thm]{Corollary}
\newtheorem{lem}[thm]{Lemma}
\newtheorem{prop}[thm]{Proposition}
\newtheorem{conj}{Conjecture}
\newtheorem{defprop}[thm]{Definition-Proposition}
\theoremstyle{definition}
\newtheorem{definition}[thm]{Definition}
\newtheorem{remark}[thm]{Remark}
\newtheorem{example}[thm]{Example}
\newtheorem{notation}[thm]{Notation}

\AtBeginEnvironment{appendices}{\crefalias{section}{appendix}}


\makeatletter
\renewcommand*\theHthm{\thesection.\arabic{thm}}
\renewcommand*\theHdefinition{\thesection.\arabic{definition}}
\renewcommand*\theHremark{\thesection.\arabic{remark}}
\renewcommand*\theHexample{\thesection.\arabic{example}}
\renewcommand*\theHprop{\thesection.\arabic{prop}}
\renewcommand*\theHlem{\thesection.\arabic{lem}}
\renewcommand*\theHcor{\thesection.\arabic{cor}}
\renewcommand*\theHnotation{\thesection.\arabic{notation}}
\makeatother




\newcommand{\centeredl}[1]{\parbox[][1.5cm][c]{3.5cm}{\centering #1}   }
\newcommand{\centered}[1]{\parbox[][1.5cm][c]{7cm}{\centering #1}   }



\newcommand{\bfp}{\boldsymbol{p}}
\newcommand{\bfx}{\boldsymbol{x}}
\newcommand{\ie}{\textit{i.e.}}


\DeclareMathOperator{\MLQ}{MLQ}
\DeclareMathOperator{\SMLQ}{MLQ^{\pm}}
\DeclareMathOperator{\Tab}{Tab}
\DeclareMathOperator{\inv}{inv}
\DeclareMathOperator{\wt}{wt}
\newcommand{\wtp}{\textup{wt}'}
\usepackage{stackengine}
\usepackage{scalerel}

\DeclareMathOperator{\skipped}{skip}
\DeclareMathOperator{\free}{free}
\DeclareMathOperator{\emp}{empty}
\DeclareMathOperator{\hook}{hook}
\DeclareMathOperator{\classic}{cl}
\DeclareMathOperator{\primed}{pr}
\DeclareMathOperator{\sign}{sign}
\DeclareMathOperator{\inc}{inc}
\DeclareMathOperator{\odd}{odd}
\DeclareMathOperator{\even}{even}
\DeclareMathOperator{\SSYT}{SSYT}


\DeclareMathOperator{\negative}{neg}
\DeclareMathOperator{\coinv}{coinv}
\DeclareMathOperator{\maj}{maj}
\DeclareMathOperator{\arm}{arm}
\newcommand{\armp}{\textup{arm}'}
\DeclareMathOperator{\leg}{leg}
\DeclareMathOperator{\ball}{ball}
\DeclareMathOperator{\pair}{pair}
\DeclareMathOperator{\rev}{rev}
\DeclareMathOperator{\Span}{Span}
\DeclareMathOperator{\Pack}{Pack}
\DeclareMathOperator{\Supp}{Supp}
\newcommand{\tH}{\widetilde{H}^{(q,t)}}
\newcommand{\xx}{\mathbf x}
\newcommand{\J}{{J}^{(\alpha)}}
\newcommand{\Jsh}{{J}^{*}}
\newcommand{\JJack}{{\mathfrak J}^{(\alpha,\gamma)}}
\newcommand{\normJ}{{j}^{(q,t)}}
\newcommand{\C}{\mathcal C}
\newcommand{\mcP}{\mathcal P}
\newcommand{\Ni}{\mathcal N^{(i)}}
\newcommand{\E}{\mathcal E}
\newcommand{\ZZ}{\mathbb Z}
\newcommand{\QQ}{\mathbb Q}
\newcommand{\QT}{\mathcal{Q}}
\newcommand{\VV}{V_{\lambda}^*}
\newcommand{\G}{{Q^{\pm}}}
\newcommand{\mcG}{\mathcal{G}}
\newcommand{\Gbar}{\widebar{\mathcal{G}}}
\newcommand{\Qbar}{\widebar{Q}}
\newcommand{\DD}{D} 
\newcommand{\NN}{\mathbb N}
\newcommand{\YY}{\mathbb Y}
\newcommand{\OO}{\mathbb{O}}
\newcommand{\On}{S_n^{\pm}}
\newcommand{\PP}{\mathbb P}
\newcommand{\MM}{T}
\newcommand{\tPP}{\widetilde{\mathbb P}}
\newcommand{\PPone}[1]{\mathbb P^{(1)}_{#1}}
\newcommand{\PPtwo}[1]{\mathbb P^{(2)}_{#1}}
\newcommand{\mfp}{\mathfrak p}
\newcommand{\mfq}{\mathfrak q}
\newcommand{\mfr}{\mathfrak r}
\newcommand{\hatfstar}[1]{\widehat{F^*_{#1}}}
\newcommand{\Estar}[2]{{E^{*#1}_{#2}}}
\newcommand{\tb}{c}

\newcommand{\fstar}[2]{{F^{*#1}_{#2}}}
\newcommand{\WS}[1]{{\widetilde{#1}}}



\usepackage[colorlinks,cite color=green,link color=purple]{hyperref}
\definecolor{green}{RGB}{43,92,47}
\definecolor{blue}{RGB}{40,68,104}
\definecolor{red}{RGB}{254, 113, 96}
\definecolor{purple}{RGB}{102,0,51}
\definecolor{gray}{RGB}{224,224,224}
\definecolor{lightpurple}{RGB}{255, 249, 242}
\definecolor{blue}{RGB}{40,68,104}
\hypersetup{colorlinks = true, urlcolor=blue}
\let\underbrace\LaTeXunderbrace


\def\LW#1{(\textcolor{teal}{LW:#1})}
\def\HBD#1{(\textcolor{violet}{HBD:#1})}


\author[]{}

%\thanks{}
%\thanks{}


\title[A probabilistic interpretation for interpolation Macdonald polynomials]{A probabilistic interpretation for\\ interpolation Macdonald polynomials}% at $q=1$}
\date{\today}
\author{Houcine Ben Dali}
\address{Department of Mathematics,
Harvard University, Cambridge, MA, and 
%}
%\address{
Center for Mathematical Sciences and Applications, Harvard University, Cambridge, MA}
\email{bendali@math.harvard.edu}
\author{Lauren Kiyomi Williams}
\address{Department of Mathematics,
Harvard University, Cambridge, MA}
\email{williams@math.harvard.edu}



\begin{document}



\maketitle

\section{The problem}

Let $\lambda=(\lambda_1 > \dots > \lambda_n \geq 0)$ be a partition with distinct parts.  Assume moreover that 
$\lambda$ is \emph{restricted}, in the sense that it has a unique part of size $0$ and no part of size $1$.
Does there exist a nontrivial Markov chain on $S_n(\lambda)$  whose stationary distribution is given by
$$\frac{F^*_{\mu}(x_1,\dots,x_n; q=1,t)}{P^*_{\lambda}(x_1,\dots,x_n;
q=1,t)} \text{ for }\mu\in S_n(\lambda)$$
where $F^*_{\mu}(x_1,\dots,x_n; q,t)$ and
$P^*_{\lambda}(x_1,\dots,x_n;q,t)$ are the
interpolation ASEP polynomial and interpolation Macdonald polynomial,
respectively?  If so, prove that the Markov chain you construct has the
desired stationary distribution.  By ``nontrivial'' we mean that the
transition probabilities of the Markov chain should not be described
using the polynomials $F_{\mu}^*(x_1,\dots,x_n; q,t)$.


\newpage 
 

\section{The solution}
\iffalse
Several recent works have given an interpretation of Macdonald polynomials in terms of stationary distributions of interacting particle models.
For example, in \cite{CantinideGierWheeler2015}, Cantini, de Gier, and Wheeler showed that the partition function of the \emph{multispecies asymmetric simple exclusion process (ASEP) on a ring} involving particles of types $\lambda_1,\dots,\lambda_n$ is equal to a Macdonald polynomial 
$P_{\lambda}(x_1,\dots,x_n;q,t)$ specialized at 
$x_1=\dots=x_n=q=1$; moreover, for each composition $\eta$ obtained by permuting the parts of $\lambda$, the  steady state probability of state $\eta$ is given by  the \emph{ASEP polynomial} $F_{\eta}$ evaluated at 
$x_1=\dots=x_n=q=1$.  (An ASEP polynomial is a certain \emph{permuted basement Macdonald polynomial}, see \cite[Proposition 4.1]{CorteelMandelshtamWilliams2022}.)

After this work, the natural question was whether one can give a parallel result which incorporates the parameters $x_1,\dots,x_n$ into the Markov chain.  This question was addressed by Ayyer, Martin, and Williams 
in   \cite{AyyerMartinWilliams2025}, who showed that the steady state probabilities and partition function of the  \emph{multispecies $t$-Push totally asymmetric simple exclusion process (TASEP)} -- a Markov chain in which there are site-dependent jump rates involving the $x_i$'s -- are given by the ASEP polynomials $F_{\eta}$ and Macdonald polynomial $P_{\lambda}$ evaluated at $q=1$.\footnote{See \cite{Defant2024} for an alternative probabilistic interpretation of these polynomials at $q=1$.}
%, a Markov chain involving particles of types $\lambda_1,\dots,\lambda_n$ hopping around a ring.  In particular, they showed that for each composition $\eta$ obtained by permuting the parts of $\lambda$, the stationary probability of being in state $\eta$ is proportional to the ASEP polynomial $F_{\eta}(x_1,\dots,x_n; 1,t)$, and the normalizing constant (or partition function) is $P_{\lambda}(x_1,\dots,x_n; 1,t)$.  

There is an inhomogeneous generalization of Macdonald polynomials due to Knop and Sahi \cite{Knop1997b, Sahi1996} called \emph{interpolation Macdonald polynomials} $P^*_{\lambda}(x_1,\dots,x_n;q,t)$, as well as an inhomogeneous generalization of ASEP polynomials called \emph{interpolation ASEP polynomials} $F^*_{\eta}(x_1,\dots,x_n;q,t)$ that we introduced in our previous work \cite{BenDaliWilliams2025}; we also gave a combinatorial formula for the interpolation ASEP and Macdonald polynomials in terms of \emph{signed multiline queues}.  
In this article 
we introduce a new Markov chain called the \emph{interpolation $t$-Push TASEP}
(see \cref{def:intpushTASEP}), which can be seen as a generalization of the $t$-Push TASEP (see \cref{rem:generalization}).  Our main theorem (see \cref{thm:main}) says that its steady state probabilities and partition function are given by the interpolation ASEP polynomials and interpolation Macdonald polynomial evaluated at $q=1$.  This generalizes the previous result of Ayyer, Martin, and Williams.




\subsection{Interpolation polynomials}
Interpolation Macdonald polynomials and interpolation ASEP polynomials are defined by vanishing conditions, as we now explain.

Given a composition $\mu=(\mu_1,\dots,\mu_n)\in\NN^n$, we define
\begin{align} \label{eq:k}
k_i(\mu)&:=\#\{j:j<i \text{ and }\mu_j>\mu_i\}+\#\{j:j>i \text{ and }\mu_j\geq \mu_i\}, \text{ and }\\\label{eq:k2}
\overline\mu&:=\left(q^{\mu_1} t^{-k_1(\mu)},\dots,q^{\mu_n} t^{-k_n(\mu)}\right).
\end{align}

%\LW{I think it would be nice for clarifying the exposition if throughout the paper, we consistently use $\lambda$ for partition and $\mu$ for composition.  Thus, in the recoloring section, e.g. \cref{prop:lumping}, we should use a different Greek letter in place of $\mu$ -- $\kappa$ might be a good choice.  (We've used $\kappa$ a tiny bit in the paper as a composition but could use another Greek letter there.) }

\emph{Interpolation Macdonald polynomials}, defined by Knop and Sahi \cite{Knop1997b, Sahi1996}, can be viewed as an inhomogeneous generalization of Macdonald polynomials.
\begin{thm}\cite{Knop1997b, Sahi1996}\label{def:intMacdonald}
For each partition $\lambda=(\lambda_1,\dots,\lambda_n)$, there is a unique inhomogeneous symmetric polynomial $P^*_{\lambda} =P_{\lambda}^*(\xx;q,t) = P_{\lambda}^*(x_1,\dots,x_n;q,t)$ 
called the \emph{interpolation Macdonald polynomial} such that 
\begin{itemize}
    \item the coefficient $[m_{\lambda}]P_{\lambda}^*$ of the 
    monomial symmetric polynomial $m_{\lambda}$ in $P_{\lambda}^*$ is $1$,
    \item $P_{\lambda}^*(\overline{\nu}) = 0$ for each partition $\nu \neq \lambda$ with $|\nu| \leq |\lambda|$.
\end{itemize}
Moreover, the top homogeneous component of 
$P_{\lambda}^*$ is the usual Macdonald polynomial
$P_{\lambda}.$
\end{thm}


 
Just as interpolation Macdonald polynomials are an inhomogeneous generalization of Macdonald polynomials, 
\emph{interpolation ASEP polynomials}, defined in \cite{BenDaliWilliams2025}, can be viewed as an inhomogeneous generalization of ASEP polynomials. 

Define $\mcP_n^{(d)}$ to be the space of polynomials in $n$ variables with degree at most $d$.

\begin{thm}[{\cite[Theorem 2.17]{BenDaliWilliams2025}}]
    Let $\lambda=(\lambda_1,\dots,\lambda_n)$ be a partition of size
    $|\lambda|=\lambda_1+ \dots + \lambda_n=d$.  For each composition $\mu\in S_n(\lambda)$, there is a unique polynomial $F^*_\mu(x_1,\dots,x_n)\in \mcP_n^{(d)}$, called the \emph{interpolation ASEP polynomial}, such that:
    \begin{itemize}
     \item for $\tau \in S_n(\lambda)$, we have 
     \begin{equation}\label{eq:normalization_F}
         [x^\tau]F^*_\mu =\delta_{\tau,\mu}.
     \end{equation}
        \item for any composition $\nu$ such that $|\nu|\leq |\mu|$ and $\nu\notin S_n(\lambda)$, we have
        $F^*_\mu(\widebar{\nu})=0$.
    \end{itemize}
Moreover, the top homogeneous component of $F^*_\mu$ is the ASEP polynomial $F_\mu$. 
\end{thm}

Just as the ASEP polynomials sum to a Macdonald polynomial, the interpolation ASEP polynomials sum to the interpolation Macdonald polynomial.  %give the symmetric Macdonald polynomials after symmetrization.

\begin{prop}[{\cite[Proposition 2.15]{BenDaliWilliams2025}}]
\label{prop:symmetrization}For any partition $\lambda$, we have $$P^*_\lambda=\sum_{\mu\in S_n(\lambda)}F^*_\mu.$$
\end{prop}

\fi
%%%%%%%%%%%%%%%%%%%%%%%%%%%%%%%%%%%%%%%%%%%%
%%%%%%%%%%%%%%%%%%%%%%%%%%%%%%%%%%%%%%%


%\subsection{The Interpolation \texorpdfstring{$t$}{t}-Push TASEP}
The answer to the question is yes, as we 
explain below.  
%We now define a Markov chain called the \emph{interpolation $t$-Push TASEP} which solves the problem.
%\begin{notation}\label{not:type}
%Let $\lambda=(\lambda_1,\dots, \lambda_n)$
%with $\lambda_1\geq\dots\geq\lambda_n\geq 0$ be a %partition. 
%We can
%describe such a partition by its vector of types
%$\mathbf{m}=(m_0, m_1, \dots, m_L)$, where
%$m_i=m_i(\lambda):=\#\{j:\lambda_j=i\}$,
%and $L$ is the largest part that occurs.
%We thus sometimes denote our partition by
%$\lambda=(L^{m_L}, \dots, 1^{m_1}, 0^{m_0})$.
%Note that $\sum_{i=0}^L m_i=n$. 
%\end{notation}
%\subsection{Notation}
 For $1\leq k\leq n$, we define
 %the following elements in $\QQ(t,x_1,\dots,x_n)$.
\begin{equation}\label{eq:1minuspk}
\mfp_k:=\frac{t^{-n+1}(1-t)}{x_k-t^{-n+2}}
\in \QQ(t,x_1,\dots,x_n) \quad \text{and}\quad \mfq_k:= \frac{(1-t)x_k}{x_k-t^{-n+2}}\in \QQ(t,x_1,\dots,x_n).
\end{equation}
%We then have
%\begin{equation}\label{eq:1minuspk}
%1-\mfp_k=\frac{x_k-t^{-n+1}}{x_k-t^{-n+2}} \qquad %\text{and}\qquad 1-\mfq_k:= \frac{t x_k-t^{-n+2}}%{x_k-t^{-n+2}}.
%\end{equation}

If  $0<t<1$ and $x_i>t^{-n+1}$ for $1\leq i\leq n$, then $\mfp_k$ and $\mfq_k$ are probabilities. 

\begin{definition}%[The interpolation $t$-Push TASEP]
\label{def:intpushTASEP}
Fix a partition $\lambda=(\lambda_1 \geq \dots \geq \lambda_n)$ with $\lambda_n=0$.
% parts labeled by $0$ are referred to as \emph{vacancies}, and parts labeled by $a>0$ are referred to as \emph{particles}. 
The \emph{interpolation $t$-Push TASEP} with 
\emph{content $\lambda$} is a Markov chain
on $S_n(\lambda)$; we think of its states as configurations of particles on a ring labeled by $\lambda_1,\dots, \lambda_n$, where state $\eta$ corresponds to having  a particle labeled $\eta_j$ at position $j$.  Moreover, there is a \emph{bell} attached to each particle.
The transitions from $\eta\in S_n(\lambda)$ are as follows.
% \LW{I changed labeling of steps to start at 0, and also correspondingly changed the labels}
%, we obtain $\nu\in S_n(\lambda)$ as follows.
\begin{enumerate}[label=(Step \arabic*)]
   \item[(Step 0)]\label{Step0} The bell at position $j$ rings with probability 
   $$P_j
   % =\frac{\prod_{k<j}\left(x_k-\frac{1}{t^{n-2}}\right)\prod_{k>j}\left(x_k-\frac{1}{t^{n-1}}\right)}{\sum_{1\leq k\leq n}\prod_{k<j}\left(x_k-\frac{1}{t^{n-2}}\right)\prod_{k>j}\left(x_k-\frac{1}{t^{n-1}}\right)}
   =\frac{\prod_{k<j}\left(x_k-\frac{1}{t^{n-2}}\right)\prod_{k>j}\left(x_k-\frac{1}{t^{n-1}}\right)}{e^*_{n-1}(\bfx;t)}, $$
   where 
   $e^*_{n-1}(\bfx;t)=\sum_{j=1}^n\prod_{k<j}\left(x_k-\frac{1}{t^{n-2}}\right)\prod_{k>j}\left(x_k-\frac{1}{t^{n-1}}\right)$.
      \item\label{Step1}%(\emph{$t$-Push TASEP}, $\circlearrowright j$)  
      The particle at position $j$, say with label $a$, is activated, and starts traveling clockwise according to the rules of the \emph{$t$-Push TASEP}.  That is, suppose  there are $m$ ``weaker'' particles in the system, i.e. particles whose labels are less than $a$,  including vacancies (label $0$).
   Then with probability $\frac{t^{k-1}}{[m]_t}$ the activated particle will move to the location of the $k$th of these weaker particles.  If this location contains a particle with positive label, then that particle becomes active, and chooses a weaker particle to displace in the same way.  The procedure continues until the active particle arrives at a vacancy.  
   %All these transitions occur simultaneously.   
   At the end of this step, position $j$ is vacant, and 
   we regard this vacancy as a particle labeled $a:=0$.

%%%%%%%%%%%%%%%%%%%%%%%%%%%%%%%%%%
\iffalse
  An equivalent way to describe the above process is that 
   each time the active particle passes a site with 
   a weaker particle, it continues to move with probability $t$, and settles at that site with probability $(1-t)$, displacing and activating the particle 
   that is located there.
   If it passes the $m$th such site, then it continues cyclically around the ring.  Note that an active particle will choose the $k$th available option with probability
   $$(1-t)(t^{k-1}+t^{k-1+m} + t^{k-1+2m}+ \dots) = 
   \frac{t^{k-1}}{1+t+\dots + t^{m-1}}.$$
   \fi
   %%%%%%%%%%%%%%%%%%%%%%%%%%%%%%%%%%%5
   %%%%%%%%%%%%%%%%%%%%%%%%%%5
   
   \item\label{Step2}%(\emph{Return to the bell}, TASEP($\rightarrow j$)) % Push TASEP"\LW{better name?}) 
   The particle labeled $a:=0$ %in position $j$ 
   now goes to position $1$ and starts traveling clockwise.
   %, looking for a new particle to activate.  
   When it gets to site $k$ for $1\leq k\leq j-1$ containing a particle with label $b \geq 0$, it 
   skips over that site with probability
   %\begin{equation}\label{eq:probskip}
    %    \begin{cases}
       $1-\mfp_k$ 
       %&\text{ 
       if $b\geq a$,
       and $1-\mfq_k$ if $b<a$;
   %\end{cases}
   %\end{equation}
      otherwise it settles at that site, activating/ displacing the site's particle.
  % with probability
   % \begin{equation}\label{eq:prob}
    %    \begin{cases}
    %   \mfp_k &\text{ if $b\geq a$},\\
    %   \mfq_k &\text{ if $b<a$}.
   %\end{cases}
   %\end{equation}
Once it activates a new particle, 
%say %with label $b\geq 0$ 
%at site $i$, 
the old particle settles at site $k$ and the new active particle continues to travel clockwise towards position $j$, activating a new particle according to the rule above. %eqref{eq:prob}.   
 %  The activated particle labeled $a\geq 0$, continues traveling clockwise, each time it finds a particle at a new position $k<j$ and labeled $b$ it activates it with probability:
  % $$\begin{cases}
   %    \mfp_k &\text{ if $b\geq a$},\\
    %   \mfq_k &\text{ if $b<a$}.
   %\end{cases}$$
   % If the particle at position $k$ is labeled $b<a$, then it is activated with probability $\mfq_k$. 
   The active particle stops once it displaces/activates another particle or arrives at position $j$, in which case it settles in position $j$.  %If a new particle is activated then it travels from left to right as above. The procedure stops when we arrive at position $j$.   
   %\LW{The particle which reaches position $j$ settles there.}
\end{enumerate}
We denote the resulting configuration by
$\nu$ and the transition probability by %$\PP_\lambda(\eta,\nu)=
$\PP(\eta,\nu)$.

Moreover, we let $\PPone{\lambda,j}=\PPone{j}$ and $\PPtwo{\lambda,j}=\PPtwo{j}$ denote the transition probabilities associated with \ref{Step1} and \ref{Step2}, respectively. We then have, for $\mu,\nu\in S_n(\lambda)$, %\LW{you didn't put a subscript $\lambda$ below.  Do you mean to drop it?}
$$\PP(\mu,\nu)=\sum_{1\leq j\leq n}P_j\sum_{\rho\in S_n(\lambda):\rho_j=0}\PPone{j}(\mu,\rho)\PPtwo{j}(\rho,\nu).$$
\end{definition}
%One can check that this Markov chain is irreducible; this implies that it has a unique stationary distribution $(\pi^*_\eta)_{\eta\in S_n(\lambda)}$.


%%%%%%%%%%%%%%%%%%%%%%%%%%
%%%%%%%%%%%%%%%%%%%%%%%%%%
\iffalse
\begin{remark}
   One may notice that the main differences between a classical step in $t$-Push TASEP (Step 1), and the new dynamics we introduce here in (Step 2) are the following: the activated particle can now displace a particle with a lower \emph{or} higher label. Moreover, the activated particle will stop with probability $1$ in position $j$; in particular, particles can no longer wrap several times around the ring.
    \end{remark}

\begin{figure}[t]
    \centering
    \includegraphics[height=0.25\linewidth]{Step_1.pdf}
    \includegraphics[height=0.26\linewidth]{Step_2.pdf}
    \includegraphics[height=0.23\linewidth]{Final_configuration.pdf}
    \caption{An example of the interpolation $t$-Push TASEP dynamics with $n=8$ and $\eta=(1,0,3,4,3,1,2,0)$. The figure on the left illustrates \ref{Step0} and \ref{Step1}: the bell rings at position $j=5$, then balls in positions 5, 7 and 1 and 2 are activated successively (the quantities next to the arrows give the transition rates). The figure in the middle illustrates \ref{Step2}: the ball at position $j=5$ is activated and travels clockwise starting from position 1. Then balls in positions 1, 2, 4<$j$ are activated successively. The figure on the right gives the final configuration.
    }
    \label{fig:Markov_chain}
\end{figure}
\begin{remark}\label{rmq:twin_particles}
Note that our assumption in \cref{def:intpushTASEP} that $\lambda$ has at least one part of size $0$ is not a restriction.  If $\lambda$ does not satisfy this property, then we can replace $\lambda$ by the partition $\tilde{\lambda} = (\lambda_1-\lambda_n, \lambda_2-\lambda_n,\dots, \lambda_{n-1}-\lambda_n, 0)$ and use the associated Markov chain.

Note also that the interpolation $t$-Push TASEP does not depend on whether a particle labeled $a$ can displace another particle labeled $a$, since they are ``identical''.
\end{remark}


\begin{example}
    We give in \cref{fig:transition_graph} the transition graph of Steps 1 and 2 in the $t$-Push$^*$ TASEP for $\lambda=(2,1,0)$, assuming that the bell rings at $j=3$ in Step 0. One may notice that, by definition, the transition probabilities in Step 2 are only defined from states $\mu$ for which $\mu_3=0$. For clarity, the loops (corresponding to transitions which do not change the configuration) are not represented in the graph. 
\end{example}


\begin{figure}[t]
    \centering
    \includegraphics[width=0.5\linewidth]{Transition_Graph.pdf}
    \caption{The transition graph of \ref{Step1} and \ref{Step2} of the $t$-Push$^*$ TASEP for $\lambda=(2,1,0)$ when $j=3$. The transition edges corresponding to \ref{Step1} (respectively \ref{Step2}) are represented in plain edges (respectively dashed edges).} 
    \label{fig:transition_graph}
\end{figure}

\begin{remark}\label{rem:generalization}
One can obtain the $t$-Push TASEP model by taking $x_i\gg t$. More formally, the asymptotics when $x_i\rightarrow \infty$ of the probabilities above give
$$P_j\sim \frac{1}{x_j},\qquad \text{and} \qquad \mfp_k\sim 0,$$
% \qquad \text{and} \qquad \mfq_k\sim 1-t.
Note in particular that the second asymptotic implies that nothing changes in \ref{Step2} (there is no particle activated, and the process stops when we arrive at position $j$).
\end{remark}
\fi
%%%%%%%%%%%%%%%%%%%%%%%%%%%
%%%%%%%%%%%%%%%%%%%%%%%%%%%%



%Our main theorem is the following.
\begin{thm}\label{thm:main}
    In the interpolation $t$-Push TASEP with content
    $\lambda=(\lambda_1,\dots,\lambda_n)$ and parameters
    $\bfx=(x_1,\dots,x_n)$ and $t$, the stationary probability of $\mu\in S_n(\lambda)$ is given by 
    $$\pi^*_{\lambda}(\mu) = \frac{F^*_{\mu}(\bfx; 1, t)}{P^*_{\lambda}(\bfx; 1, t)}.$$
    %where $F^*_{\mu}(\bfx;q,t)$ is the interpolation ASEP polynomial, and $P^*_{\lambda}(\bfx; q, t)$ is the interpolation Macdonald polynomial.
\end{thm}


%%%%%%%%%%%%%%%%%%%%%%%%%%%%%%%%%%%%%%%%%
%%%%%%%%%%%%%%%%%%%%%%%%%%%%%%%%%%%%%%%%%%%%%5
\iffalse
A combinatorial formula for the interpolation ASEP polynomials in terms of \emph{signed multiline queues} was given in \cite[Theorem 1.15]{BenDaliWilliams2025}.  Combining this with \cref{thm:main}, we obtain the following corollary.

\begin{cor}
Fix a partition $\lambda=(\lambda_1,\dots,\lambda_n)$, and consider the signed multiline queues of type $\mu$, for $\mu\in S_n(\lambda),$ with parameters $x_1,\dots,x_n,t$ and $q=1$.
The distribution of the bottom line of the signed multiline queues is the same as the stationary distribution of the interpolation $t$-Push TASEP with content $\lambda$.
\end{cor}

\subsection{Strategy of Proof}
Let us briefly describe the strategy for proving \cref{thm:main}, which is similar to the strategy used in 
\cite{AyyerMartinWilliams2025} for proving the analogous result connecting the $t$-Push TASEP to ASEP polynomials.
We start by giving a combinatorial proof of the theorem
when $\lambda$ has distinct parts.  We use the fact that 
interpolation ASEP polynomials are generating functions for 
signed multiline queues \cite{BenDaliWilliams2025}, 
and that signed multiline queues can be built by stacking alternating layers of ``classic'' and ``signed'' layers of pairings of balls.  We also use the fact that
$\PPone{\lambda,j}(\eta, \kappa)$ gives the probability of obtaining
$\kappa$ as Row $1'$ of a signed multiline queue with content $\lambda$, given that Row $2$ is $\eta$ and the vacancy in Row $1'$ is at site $j$.  Similarly, we show that 
$\PPtwo{\lambda,j}(\kappa, \nu)$ gives the probability of obtaining $\nu$ as Row $1$ of a signed multiline queue with content $\lambda$, given that Row $1'$ is $\kappa$ and the vacancy in Row $1'$ is at site $j$.  This allows us to directly connect the transition probabilities of our Markov chain to the generating functions for signed multiline queues, which give the interpolation ASEP polynomials.

Having proved \cref{thm:main} for partitions with distinct parts, we generalize the result to all partitions, by 
using the fact that ``recoloring'' our particles in the Markov chain via a weakly order-preserving function $\phi:\NN \to \NN$ gives rise to a projection or lumping from a ``finer'' to a ``coarser'' version of the interpolation $t$-Push TASEP.  This allows us to compute the stationary distribution of a ``coarser'' interpolation $t$-Push TASEP associated to a partition with repeated parts, using the stationary distribution coming from a partition with distinct parts.  Meanwhile, we prove that interpolation ASEP polynomials at $q=1$ satisfy a corresponding identity that gives a compatibility relation when one applies a weakly order-preserving function.



\bigskip
% \HBD{I added the following subsection.} \LW{great!}
\subsection{On the multispecies ASEP}
As mentioned earlier,   Cantini, de Gier, and Wheeler \cite{CantinideGierWheeler2015} showed that
the specialization $P_{\lambda}(x_1=\dots=x_n=1;q=1,t)$ of the homogeneous Macdonald polynomial 
 is the partition function of the multispecies ASEP on a ring. This result follows readily from the fact that the homogeneous ASEP polynomials $F_\mu$ are characterized by a family of equations called the \emph{qKZ equations} (see e.g. \cite[Equations (17) to (19)]{CantinideGierWheeler2015}). 

A natural question is to find an interpolation analogue of the multispecies ASEP, in the same spirit as the analogue we give here for the $t$-Push TASEP model. A main obstacle to such a generalization is the fact that interpolation ASEP polynomials $F_\mu^*$ do not satisfy all the qKZ equations (more precisely, they lack the circular symmetry property). It would then be interesting to find an algebraic characterization of the interpolation polynomials $F_\mu^*$ from which the qKZ equations can be recovered by taking the top homogeneous part.

\HBD{Should we also mention (as another obstacle) that the interpolation polynomials are in general not positive at $x_1=\dots=x_n$?}

\medskip 

The structure of this paper is as follows. In \cref{sec:MLQ}, we review the notions of  two-line queues and signed two-line queues. \cref{sec:recoloring_Push_TASEP} is devoted to a recoloring property of the interpolation $t$-Push TASEP model. In \cref{sec:interpolation_polynomials}, we establish several properties of the interpolation ASEP polynomials at $q=1$.  We then prove the main theorem in \cref{sec:proof}. In \cref{sec:density}, we derive the density formula. Finally, in \cref{app:shape_permuting,app:permuted_basement,app:t_int_Schur}, we collect useful properties of interpolation ASEP polynomials, along with a Jacobi--Trudi identity for symmetric interpolation polynomials at $q=t$.



\bigskip
\noindent{\bf Acknowledgements:~}
HBD acknowledges support from the Center of Mathematical Sciences and Applications at Harvard University.
LW was supported by the National Science Foundation under Award No.
DMS-2152991.
%until May 12, 2025, when the grant was terminated; she would also like to thank the Radcliffe Institute for Advanced Study, where some of this work was carried out. 
Any opinions, findings, and conclusions or recommendations expressed in this material are
those of the author(s) and do not necessarily reflect the views of the National Science
Foundation.  


\section{Multiline queues}\label{sec:MLQ}

\emph{Muliline queues} are a combinatorial object whose generating functions
encode Macdonald polynomials \cite{CorteelMandelshtamWilliams2022}; roughly speaking, multiline queues can be built by stacking \emph{two-line queues}.
\emph{Signed multiline queues}
generalize multiline queues, and are a combinatorial object whose generating functions encode interpolation ASEP polynomials \cite{BenDaliWilliams2025}; roughly speaking, signed multiline queues 
can be built by stacking alternating sequences of (classical) two-line queues and \emph{signed two-line queues}.  
In this section we define both 
(classical) two-line queues and signed two-line queues.  We will later explain how these two-line queues encode transitions in the 
interpolation $t$-Push TASEP.

\fi
%%%%%%%%%%%%%%%%%%%%%%%%%%%%%%%%%%%%%%%%%5
%%%%%%%%%%%%%%%%%%%%%%%%%%%%%%%%%%%%%%%

%%%%%%%%%%%%%%%%%%%%%%%%%%%%%%%%%%%%5
%%%%%%%%%%%%%%%%%%%%%%%%%%%%%%%%%
\iffalse 
\subsection{(Classical) two-line queues}

\begin{definition}\label{def:paired_system}
A \emph{paired ball system} is a $2 \times n$ array
in which each position is either occupied by a ball or is empty,
and there are at least as many balls on the bottom row as on the top;
moreover, each ball in the top row is paired with one in the bottom, and joined by a shortest strand traveling weakly left to right, allowing wrapping if necessary.
%    A \emph{paired ball system} consists of a ball system
 %   on two rows, with the same number of balls on each row, in which balls are paired (by a shortest strand traveling weakly left to right, allowing wrapping if necessary), and paired balls have the same positive integer label of $2$ or more.
  %Empty positions are sometimes represented by a ball labeled $0$.
\end{definition}


\begin{definition}\label{def:two-row}
A \emph{two-line queue}\footnote{In \cite{CorteelMandelshtamWilliams2022,BenDaliWilliams2025}, these objects are referred to as generalized two-line queues. For simplicity, we will call them two-line queues throughout this paper.} is a paired ball system,   such that:
\begin{enumerate}
\item each ball is labeled by an element of $\ZZ^+$;
\item paired balls have the same label;
\item\label{GMLQ:item1} if a ball with label $a\in \ZZ^+$ in the top row has a ball directly underneath it labeled $a'$, then
        we must have $a'\geq a$, and if $a'=a$, the two balls must be paired to each other. %trivially paired.
       % \item \LW{do we want to require that any singleton ball on the 
        %bottom row has label $1$?}
\end{enumerate}
The bottom and top rows of a two-line queue are then represented by  compositions $\mu$ and $\nu$ in $\NN^n$. \HBD{where an empty position is represented by 0.}  We refer to a pairing of two balls in the same column as a \emph{trivial pairing}.
%Let $\mcG^\alpha_\mu$ denote the
%set of (generalized) signed two-line queues %with bottom row $\mu$ and top row $\alpha$. 
\end{definition} 


\begin{figure}[t]
    \centering
    \includegraphics[width=0.5\linewidth]{2_line_queue.pdf}
    \caption{An example of a two-line queue in $\mathcal{Q}^{(4,0,2,0,0,3,0,0,4)}_{(0,4,4,0,2,3,0,1,0)}.$}
    \label{fig:two_line_queue}
\end{figure}

See \Cref{fig:two_line_queue} for an example of a two-line queue.

\begin{remark}\label{rem:MLQ}
    More generally, a \emph{multiline queue}
    is a vertical concatenation of 
    two-line queues, with rows labeled from $1$ to $L$ from bottom to top.
    We refer to a collection of linked pairs of balls as a \emph{strand},
    and we require that all balls in a strand whose top ball is in row $r$, are labeled $r$.
    See \cite{CorteelMandelshtamWilliams2022} for the precise definition.
   %    can be viewed as the bottom two rows in a usual multiline queue.
%Similarly for the signed two-line queues.}
\end{remark}


We associate to each pairing $p$ in a two-line queue $Q$ a weight $\wt(p)$ defined as follows: we read the balls in the top row in decreasing order of  their label; within a fixed  value, we read the balls from right to left.  Reading the balls in this order, we place the strands pairing the balls one by one.  The balls in the bottom row that have not yet been matched  are \emph{free}.  If a pairing $p$ matches a ball labeled $a$ in the top row and column $j$ to a ball labeled $a$ in the bottom row and column $k$, then the free balls in the bottom row and columns $j+1,j+2,\dots,k-1$ (where this set of columns is cyclically consecutive)
are  \emph{skipped}.  We then set 
\begin{equation}
\wt(p) = 
         (1-t) \frac{t^{\skipped(p)}}{1-t^{\free(p)}} 
\end{equation}
Having associated a weight to each nontrivial pairing,
we define $$\wt_{\pair}(Q) = \prod_p \wt(p),$$
where the product is over all nontrivial pairings of balls in $Q$.

We let  $\mathcal{Q}_\kappa^\eta$ denote the set of 
classical two-line queues with top row $\eta$ and bottom row $\kappa$, and we let
$a^\eta_{\kappa}$  
denote  the weight generating function of these two-line queues:\footnote{We note that in \cite{CorteelMandelshtamWilliams2022}, 
two-line queues were defined together with a weight function that depends on both $t$ and $q$; here we specialize to $q=1$.}
\begin{equation}\label{eq:generating_function0}
    a_\kappa^\eta=a_\kappa^\eta(t):=\sum_{Q\in\mathcal{Q}_\kappa^\eta}\wt_{\pair}(Q).
\end{equation}
\fi
%%%%%%%%%%%%%%%%%%%%%%%%%%%%%%%%%5
%%%%%%%%%%%%%%%%%%%%%%%%%%%%%%%%%%%%%%%%

\newpage 

%\subsection{Signed two-line queues}
\section{The proof} 

Recall the notion of classical two-line queues 
from \cite{CorteelMandelshtamWilliams2022} and signed two-line queues from \cite{BenDaliWilliams2025} together with their weight functions.  (Here we specialize  $q=1$.)
%to give a combinatorial formula for the interpolation ASEP polynomials $F^*_\mu(\bfx;q,t)$. 



Let  $\mathcal{Q}_\kappa^\eta$ denote the set of 
classical two-line queues with top row $\eta=(\eta_1,\dots,\eta_n)$ and bottom row $\kappa=(\kappa_1,\dots,\kappa_n)$, and  let
$a^\eta_{\kappa}$  
denote  the weight generating function of $\mathcal{Q}_{\kappa}^{\eta}$.
\begin{equation}\label{eq:generating_function0}
    a_\kappa^\eta=a_\kappa^\eta(t):=\sum_{Q\in\mathcal{Q}_\kappa^\eta}\wt_{\pair}(Q).
\end{equation}

Let $\mcG^\alpha_\mu$ denote the
set of signed two-line queues with 
top row $\alpha=(\alpha_1,\dots,\alpha_n)$ and 
bottom row $\mu=(\mu_1,\dots,\mu_n)$, and let
$b_\mu^\alpha$  denote the weight generating function of $\mcG^\alpha_\mu$. 
\begin{equation}\label{eq:generating_function1}
    b_\mu^\alpha=b_\mu^\alpha(t):=\sum_{Q\in\mcG_\mu^\alpha}\wt_{\pair}(Q).
\end{equation}
%%%%%%%%%%%%%%%%%%%%%%%%%
%%%%%%%%%%%%%%%%%
\iffalse 
We give here the case $q=1$ of this result.

\begin{definition}[\cite{BenDaliWilliams2025}]\label{def:signed_two-row}
A \emph{signed two-line queue} is a paired ball system (as in \cref{def:paired_system}), 
 such that:
\begin{enumerate}
 \item\label{GSMLQ:item1} Each pairing connects two balls with a shortest strand that travels either straight down or from left to right, and does not wrap around;
\item each ball in the top row is labeled by an element of $\ZZ$ \HBD{$\ZZ\backslash\{0\}?$} and each ball in the bottom row is labeled by an element of $\ZZ^+$;
\item paired balls have labels with the same absolute value;
 \item\label{GSMLQ:item2} a ball with label $a\in \ZZ^+$ in the top row
            must always have a ball labeled $a'$ underneath it, where $a' \geq a$, and if $a'=a$, the two balls must be trivially paired;
    \item\label{GSMLQ:item3} In the top row, each negative ball with label $-a$ (for $a\in \ZZ^+$) has either an empty spot below it or a ball with label $a'$, where $a \geq a'$.
\end{enumerate}
\end{definition} 
The balls in the bottom row of a signed two-line queue are called \emph{regular balls}, and those in the top row are called \emph{signed balls}.
The bottom row is then represented by a composition $\mu\in \NN^n$, and the top row by a signed permutation $\alpha$ of $\mu$. Let $\mcG^\alpha_\mu$ denote the
set of signed two-line queues with bottom row $\mu$ and top row $\alpha$. 

See \Cref{fig:ghost_two_line_queue} for an example of a signed two-line queue.


\begin{remark}\label{rem:sMLQ}
As in \cref{rem:MLQ}, a \emph{signed multiline queue}
    is a vertical concatenation of 
    two-line queues, which alternate between signed and (usual) two-line queues, with rows labeled $1, 1', 2, 2', \dots, L, L'$ from bottom to top.
    We refer to a collection of linked pairs of balls as a \emph{strand},
    and we require that all balls in a strand whose top ball is in row $r$, have labeled which is $\pm r$.
    See \cite{BenDaliWilliams2025} for the precise definition.
   %    can be viewed as the bottom two rows in a usual multiline queue.
%Similarly for the signed two-line queues.}
\end{remark}


\begin{figure}[t]
    \centering
    \includegraphics[width=0.5\linewidth]{signed_2_line_queue.pdf}
    \caption{An example of a signed two-line queue in $\mcG_{(0,0,4,4,0,0,3,2,1)}^{(0,4,-4,0,-2,-3,0,1,0)}.$
    }
    \label{fig:ghost_two_line_queue}
\end{figure}



%\subsection{Weights and generating function}
% Using \eqref{eq:pair2}, we define the weight of a signed two-line queue $\G\in \mcG^\alpha_\mu$ to be
% $$\wt_{\pair}(\G)=\prod_{p}\wt_{\pair}(p),$$ where the product is over all nontrivial (signed) pairings of $\G$. We then define the weight generating function $G_\mu^\alpha$ of $\mcG^\alpha_\mu$ to be $$G_\mu^\alpha=G_\mu^\alpha(t):=\sum_{\G\in\mcG_\mu^\alpha}\wt_{\pair}(\G).$$



% Recall that $\skipped$ is the statistic associated to a nontrivial pairing defined in \cref{def:wt}. We now give an equivalent ``static'' description of this statistic which follows directly from the definitions.
% \begin{lem}\label{lem:skipped}
%     Let $\G$ be a signed two-line queue and let $p$ be a nontrivial $a$-pairing connecting a signed ball labeled $\pm a$ in column $i$ of the top row, to a regular ball in column $j>i$ in the bottom row labeled $a$. Then $\skipped(p)$ counts the number of balls $B$ in the bottom row and in a column $r$ labeled by $c\in \NN_+$, such that $i<r<j$ and, either
%     \begin{itemize}
%         \item $c<a$,
%         \item or $c=a$ and the ball to which $B$ is paired lies in a column $k<i$.
%     \end{itemize}
% \end{lem}





We associate to each pairing $p$ in a signed two-line queue a weight $\wt(p)$ defined as follows: we read the balls in the top row in decreasing order of the absolute value of their label; within a fixed absolute value, we read the balls from right to left.  Reading the balls in this order, we place the strands pairing the balls one by one.  The balls in the bottom row that have not yet been matched  are \emph{free}.  If a pairing $p$ matches a ball labeled $\pm a$ in the top row and column $j$ to a ball labeled $a$ in the bottom row and column $k>j$, then the free balls (respectively, empty positions) in the bottom row and columns $j+1,j+2,\dots,k-1$
are  \emph{skipped} (respectively, \emph{empty}).  We then set 
\begin{equation}\label{eq:pair2}
\wt(p) = \begin{cases}
         (1-t) t^{\skipped(p)+\emp(p)} &\text{ if $p$ connects a positive ball and a regular ball}\\
         -(1-t) t^{\skipped(p)+\emp(p)} &\text{ if $p$ connects a negative ball and a regular ball.}
\end{cases}
\end{equation}
Having associated a weight to each nontrivial pairing,
we define $$\wt_{\pair}(Q) = \prod_p \wt(p),$$
where the product is over all nontrivial pairings of balls in $Q$.
We define the weight generating function $b_\mu^\alpha$ of $\mcG^\alpha_\mu$ to be
\begin{equation}\label{eq:generating_function1}
    b_\mu^\alpha=b_\mu^\alpha(t):=\sum_{Q\in\mcG_\mu^\alpha}\wt_{\pair}(Q).
\end{equation}

\begin{remark}
In this paper, we are interested in the case when $\lambda$ has only one vacancy i.e $m_0(\lambda)=1$. In this case, one can check that for any pairing $p$ we have $\emp(p)=0$, and that $\skipped(p)$ has the following simpler description: if  $p$ matches a ball labeled $\pm a$ in the top row and column $j$ to a ball labeled $a$ in the bottom row and column $k>j$, then $\skipped(p)$ corresponds to the number of balls in the bottom row and in a column $j+1,j+2,\dots,k-1$ which are labeled $b<a$.
    
\end{remark}



We assign to each ball $B$ in the top row in column $k$ of a signed two-line queue $Q$ a weight
$$\wt(B):=\begin{cases}
    x_k &\text{if $B$ is positive}\\
    -\frac{1}{t^{n-1}}&\text{if $B$ is negative}.
\end{cases}$$

The total weight of $Q$ is then defined by 
$$\wt(Q):=\wt_{\pair}(Q)\prod_{B \text{ in the top row}}\wt(B).$$

Adding the ball weights to \cref{eq:generating_function1} we obtain 
\fi
%%%%%%%%%%%%%%%%%%%%%%%%%%%%%%%%%5
%%%%%%%%%%%%%%%%%%%%%%%%%%%
Let $\wt(Q):=\wt_{\pair}(Q)\wt_{\ball}(Q)$ be the product of the pair weight and \emph{the ball weight}.

We obtain
\begin{equation}\label{eq:generating_function2}
    \wt_\alpha b_\mu^\alpha=\sum_{Q\in\mcG_\mu^\alpha}\wt(Q),
    \ \text{ where } \ 
  \wt_\alpha:=\prod_{k:\,\alpha_k>0}x_k  \prod_{k:\, \alpha_k<0}\frac{-1}{t^{n-1}}.
\end{equation}


%\subsection{Unsigned versions of multiline queues}
%\label{ssec:MLQ_distinct_parts}
%In this section we assume that all parts of $\lambda$ are distinct ($m_i(\lambda)\leq 1$). We give an unsigned version of multiline queues in this case, which we use to prove that the generating functions $\wt_\alpha b_\mu^\alpha$ is actually positive under some assumptions on the parameters $t, x_1,\dots, x_n$.
%This property will be useful to obtain a Markov chain with positive transition rates.


\begin{definition}\label{def:Gbar}
Given a signed two-line queue $Q\in \mcG^\alpha_{\mu}$, we associate to it an \emph{unsigned version} $\Qbar$ obtained by forgetting the signs of the balls in the top row. The composition we read in the bottom row (respectively the top row) of $\Qbar$ is $\mu$ (respectively $\lVert \alpha\rVert)$, where 
$$\lVert \alpha\rVert=(|\alpha_1|,\dots,|\alpha_n|).$$
We then define $\Gbar_\mu^{\kappa}$ as the set of paired ball systems obtained by applying this operation on $Q\in \mcG^\alpha_\mu$, where $\alpha\in \ZZ^n$ satisfying $\lVert\alpha\rVert=\kappa$.
\end{definition}

This leads us to define the following weights. Fix $\Qbar\in \Gbar^{\kappa}_\mu$:
\begin{itemize}
    \item A nontrivial pairing $p$ in $\Qbar$ has the weight
    \begin{equation}\label{eq:pairing_weight_unsigned}
  \wt(p)=(1-t)t^{\skipped(p)}.      
    \end{equation}

    \item Let $B$ be a ball labeled $a>0$ in column $k$ and such that the ball below is labeled $b$ (If $B$ has a vacancy below it, we take $b=0$.) We define the weight of $B$ by:
\begin{equation}\label{eq:weights_unsigned_balls}
\wt(B):=
\begin{cases}
    x_k-\frac{1}{t^{n-1}}&\text{if $b=a$,}\\
    x_k&\text{if $b>a$,}\\
    \frac{1}{t^{n-1}}&\text{if $b<a$}.
\end{cases}    
\end{equation}    
The weight of $\Qbar$ is defined by
$$\wt(\Qbar):=\prod_{B \text{ in the top row}}\wt(B) \prod_{p \text{ nontrivial pairing}}\wt(p).$$
\end{itemize}
 We then have the following lemma.


 
 \begin{lem}\label{lem:forget_signs}
 Fix a partition $\lambda$ with distinct parts and two compositions $\kappa,\mu\in S_n(\lambda)$. Let $\Qbar\in \Gbar^{\kappa}_\mu$. Then
 $$\wt(\Qbar)=\sum_{Q}\wt(Q),$$
where the sum is taken over all signed two-line queues $Q$ from which $\Qbar$ is obtained by forgetting signs.
 \end{lem}
\begin{proof}
    We consider all the possible ways of ``adding signs'' to the balls in the top row of $\Qbar$ to obtain a signed two-line queue. Fix such a ball $B$ labeled $a>0$:
    \begin{itemize}
        \item if $B$ has below it a vacancy or a ball labeled $b<a$, then 
        %from \cref{def:signed_two-row} \cref{GSMLQ:item2}, w
        we must assign a $-$ sign to $B$. 
        \item if $B$ has a ball labeled $b>a$ below it, then 
        %\cref{def:signed_two-row} \cref{GSMLQ:item3}, 
        we must assign a $+$ sign to $B$.
        \item if $B$ has a ball labeled $b=a$ below it, then %from 
        %\cref{def:signed_two-row} \cref{GSMLQ:item1} and the fact that $\lambda$ has distinct parts we know that the two balls should be trivially paired. In this case 
        we can give $B$ a $+$ or $-$ sign.
    \end{itemize}
We then check that the possible signs for each ball $B$ is consistent with the choice of weights in \cref{eq:weights_unsigned_balls}. In particular, one notices that when a ball $B$ is given a $-$ sign, the ball weight should be multiplied by $-1$ when we go from $\Qbar$ to $Q$, but the weight of the pairing connected to $B$ is also multiplied by $-1$.
%(compare \cref{eq:pair2,eq:pairing_weight_unsigned}).
\end{proof}

Given $\kappa\in S_n(\nu)$, we define $\tb_\nu^{\kappa}$ by 
\begin{equation}\label{eq:def_c}
  \tb_{\nu}^{\kappa}:=\sum_{\alpha:\, \lVert \alpha\rVert =\kappa}\wt_\alpha b^{\alpha}_{\nu}.  
\end{equation}


 We get the following corollary obtained by combining \cref{eq:generating_function2} and \cref{lem:forget_signs}.
\begin{lem}\label{lem:c}
 Fix $\lambda$ a partition with distinct parts, and $\kappa,\mu\in S_n(\lambda)$. Then
 $$c^\kappa_\mu=\sum_{\Qbar\in \Gbar^{\kappa}_\mu} \wt(\Qbar).$$
\end{lem}
    Since $\lambda$ has distinct parts, $\Gbar^{\kappa}_\nu$ is either empty or contains exactly one element.



%%%%%%%%%%%%%%%%%%%%%%%%%%%
\iffalse 
\begin{figure}
    \centering
    \includegraphics[width=0.5\linewidth]{UGTLQ.pdf}
    \caption{The unique paired ball system in $\Gbar^{(2,7,1,5,0,6)}_{(0,7,2,1,5,6)}.$}
    \label{fig:UGTLQ}
\end{figure}
\fi
%%%%%%%%%%%%%%%%%%%%%%


 
%The following result is from \cite[Theorem 1.15 and Lemma 5.6]{BenDaliWilliams2025}.
%     \begin{thm}[{\cite{BenDaliWilliams2025}}]\label{lem:F_decomposition}
 %For any  composition $\nu\in \NN^n$, we have
 %\begin{equation}\label{eq:F_decomposition}
 %  F^*_\nu(\bfx;1,t)=\sum_{\eta\in\NN^n}{F^{*\eta}_\nu}(\bfx;1,t) F^*_{\eta^{-}}(\bfx;1,t),  
 %\end{equation}
 %where 
 %$$F^{*\eta}_{\nu} :=\sum_{\alpha\in\ZZ^n}b_\nu^\alpha \wt_\alpha a^\eta_{\lVert \alpha\rVert}=\sum_{\mu\in \NN^n}a_\mu^\eta c^\mu_\nu .$$
 %\end{thm}



\begin{comment}
\subsection{Some useful computation}
All the coefficients we consider in this section are evaluated at $q=1$.
\begin{lem}
    Fix a signed composition $\alpha\in \ZZ^n$, with only one vacancy at position $i$, and fix $1\leq j\leq n.$ Then
\begin{equation}\label{eq:fixed_j}
    \sum_{\mu:\mu_j=0}b^\alpha_\mu=\begin{cases}
        0 &\text{ if }j>i\text{ or }\alpha_j>0,\\
        1 &\text{ if }j=i,\\
        -(1-t)t^{\#\left\{k:\ j<k<i \text{ and }\alpha_k<0 \right\}}&\text{ if }j<i\text{ and }\alpha_j<0.
    \end{cases}
\end{equation}
\end{lem}

For $\alpha\in \ZZ^n$, we define 
$\sign(\alpha)\in\{+1,-1,0\}$ as the sequence of signs of $
\alpha$. By summing \cref{eq:fixed_j} over all $j$, we deduce the following.
\begin{cor}
Fix a signed composition $\alpha\in \ZZ^n$ such that there is a unique $i$ with $\alpha_i=0$. %$, with only one vacancy at position $i$. 
Then
    $$\sum_{\mu}b^\alpha_\mu=t^{\#\left\{k: k<i \text{ and }\alpha_k<0 \right\}.}$$
    In particular, this sum  depends only on $\sign(\alpha).$
\end{cor}


Recall that 
$$\fstar{}{\mu}=\sum_{\eta}\fstar{\eta}{\mu}F^*_\eta,$$
with
$$\fstar{\eta}{\mu}=\sum_{\alpha\in \ZZ^n}b^\alpha_{\mu}\wt_\alpha a^\eta_{\lVert\alpha\rVert}.$$

% Recall \LW{If this is a "recall" from our other paper, better to leave out the word} also that, 
If $\alpha\in \ZZ^n$ with $\sign(\alpha)=\varepsilon$, we define
$$\wt_\alpha=\wt_{\varepsilon}:=\prod_{k:\ \alpha_k>+1}x_k\prod_{k:\ \alpha_k<0}\frac{-1}{t^{n-1}}.$$

We start from a configuration $\eta$, with one vacancy at position $i$.
Fix a sequence $\varepsilon=(\varepsilon_1,\dots,\varepsilon_{j-1},0,\varepsilon_{j+1},\dots,\varepsilon_n)$ with $\varepsilon_k\in\{\pm 1\}$ for $k\neq j.$
We want that the probability of the bell to ring at position $j$, and that $\sign(\alpha)=\varepsilon$ is proportional to
  
\begin{align*}
\wtp(\varepsilon):=\sum_{\substack{\alpha\in \ZZ^n\\ \sign(\alpha)=\varepsilon}}\sum_{\mu}b^\alpha_{\mu}\wt_\alpha a^\eta_{\alpha}
&=\wt_\varepsilon 
\sum_{\substack{\alpha\in \ZZ^n\\ \sign(\alpha)=\varepsilon}}a^\eta_{\lVert\alpha\rVert}
\sum_{\mu}b^\alpha_{\mu}\\
&=\wt_\varepsilon t^{\#\left\{k:\ k<i \text{ and }\varepsilon_k=-1 \right\}}
\sum_{\substack{\nu\in \NN^n\\ \nu_k=0\Leftrightarrow k=j}}a^\eta_{\nu}\\
&=\wt_\varepsilon t^{\#\left\{k:\ k<i \text{ and }\varepsilon_k=-1 \right\}}\\
&=\prod_{k:\ \varepsilon_k=+1}x_k\prod_{k<j:\ \varepsilon_k=-1}\frac{-t}{t^{n-1}}\prod_{k>j:\ \varepsilon_k=-1}\frac{-1}{t^{n-1}}.
\end{align*}
\end{comment}




%%%%%%%%%%%%%%%%%%%%%%%%%%%%%%%%%%%%%%%%%%%%%%55
%%%%%%%%%%%%%%%%%%%%%%%%%%%%%%%%%%%%%%%%%%%%
\iffalse 
\section{Recoloring the interpolation \texorpdfstring{$t$}{t}-Push TASEP}\label{sec:recoloring_Push_TASEP}
Following \cite{AyyerMartinWilliams2025}, we introduce the following definition.
\begin{definition}
    We say that a function $\phi$ from $\NN$ to $\NN$ is \emph{weakly order-preserving} if $\phi(i)\leq \phi(j)$ whenever $i\leq j$.
    % , and such that\footnote{This second condition is not in the definition of \cite{AyyerMartinWilliams2025}.} $\phi(0)=0$. If $\phi$ is such a function, we extend it to a function from $\OO$ to $\OO$ in such a way that $\phi(0^*)=0^*$, $\phi(+i)=+\phi(i)$ and $\phi(-i)=-\phi(i)$, for any $i\in \NN$.
\end{definition}
If $\mu\in \NN^n$ is a composition, we define $\phi(\mu):=(\phi(\mu_1),\dots,\phi(\mu_n))\in \NN^n$.
% Similarly, if $\alpha\in \OO^n$, then $\phi(\alpha):=(\phi(\alpha_1),\dots,\phi(\alpha_n))\in \OO^n$. 
Note that if $\mu$ is a partition then so is $\phi(\mu)$. In what follows, we will often apply a weakly order-preserving map to states of the interpolation $t$-Push TASEP; we refer to this as 
\emph{recoloring}.


\begin{prop}\label{prop:lumping}
    Let $\phi:\NN \to \NN$ be a weakly order-preserving function with $\phi(0)=0$, and let $\lambda=(\lambda_1,\dots,\lambda_n)$ and $\kappa=(\kappa_1,\dots,\kappa_n)$ be  partitions such that $\kappa=\phi(\lambda)$.
Then the interpolation $t$-Push TASEP with particle content $\lambda$ lumps (or projects) to the interpolation $t$-Push TASEP with particle content $\kappa$.  That is,  
    if $\eta,\zeta\in S_n(\kappa)$ and $\mu\in S_n(\lambda)$ such that $\phi(\mu)=\eta$, then
    $$\sum_{\nu\in S_n(\lambda):\, \phi(\nu)=\zeta}\PP_{\lambda}(\mu,\nu)=\PP_\kappa(\eta,\zeta).$$
\end{prop}

\begin{proof}
We note that each of the three quantities $(P_j,\PPone{},\PPtwo{})$ defining the transition rates in the interpolation $t$-Push TASEP behaves well under recoloring: $P_j$ is independent from the state $\mu$, and for any $j$, we have:
\begin{align*}
  &\sum_{\rho\in S_n(\lambda):\, \phi(\rho)=\theta}\PPone{\lambda,j}(\mu,\rho)
  =\PPone{\kappa,j}(\eta,\theta),\quad \text{ for any $\eta,\theta\in S_n(\kappa)$ and $\mu\in S_n(\lambda)$ with $\phi(\mu)=\eta$,}\\
    &\sum_{\nu\in S_n(\lambda):\, \phi(\nu)
    =\zeta}\PPtwo{\lambda,j}(\rho,\nu)=\PPone{\kappa,j}(\theta,\zeta),\quad \text{ for any $\theta,\zeta\in S_n(\kappa)$ and $\rho\in S_n(\lambda)$  with $\phi(\rho)=\theta$.}
\end{align*}

 %Thus, as in \cite[Proposition 2.2]
 %{AyyerMartinWilliams2025}, 
 \cref{prop:lumping} follows from the above observations, %description of the dynamics, see 
 \cref{rmq:twin_particles}, and the fact that the 
 state space of the first Markov chain projects onto the state space of the second.
\end{proof}


\begin{figure}[t]
    \centering
    \includegraphics[height=0.25\linewidth]{Step_1_restricted.pdf}
    \includegraphics[height=0.26\linewidth]{Step_2_restricted.pdf}
    \includegraphics[height=0.23\linewidth]{Final_configuration_restricted.pdf}
    \caption{An example of the interpolation $t$-Push TASEP dynamics with $n=8$ for a restricted partition $\lambda=(8,7,6,5,4,3,2,0)$ and a state $\mu=(3,0,6,8,7,4,5,2)\in S_n(\lambda)$.}
    % The figure on the left illustrates \ref{Step0} and \ref{Step1}: the bell rings at position $j=5$, then balls in positions 5, 7 and 1 and 2 are activated successively (the quantities next to the arrows give the transition rates). The figure in the middle illustrates \ref{Step2}: the ball at position $j=5$ is activated and travels clockwise starting from position 1. Then balls in positions 1, 2, 4<$j$ are activated successively. The figure on the right gives the final configuration.
    % \LW{We'll have a similar figure but with ball content
    % $(0, 2,3,4,5,6,7,8)$ in order to illustrate both the projection property
    % and also the connection with two-line queues.}}
    \label{fig:Markov_chain_2}
\end{figure}

\begin{example}
    In \cref{fig:Markov_chain_2}, we illustrate a transition for the $t$-Push$^*$ TASEP indexed by $\lambda=(8,7,6,5,4,3,2,0)$.  The corresponding transition for the $t$-Push$^*$ TASEP indexed by indexed by $\kappa=(4,3,3,2,1,1,0,0)$ from \cref{fig:Markov_chain} is then obtained by projecting, using for example by the map:
    $$\phi(i)=\begin{cases}
        0 & \text{if $i\in \{0,1,2\}$}\\
        1 & \text{if $i\in \{3,4\}$}\\
        2 & \text{if $i\in \{5\}$}\\
        3 & \text{if $i\in \{6,7\}$}\\
        4 & \text{if $i\geq 8$.}\\
    \end{cases}$$
\end{example}

We then have the following consequence on the stationary distribution. It follows from the strong lumping property of Markov chains (see \textit{e.g} \cite{KemenySnell1976,Pang2019,Williams2022}.)

\begin{cor}\label{cor:lumping}
Let $\phi,\lambda$ and $\kappa$ as above. Then  the stationary distributions of the 
interpolation $t$-Push TASEPs with particle contents $\kappa$ and $\lambda$ are related as follows:
for any $\eta\in S_n(\kappa)$, 
$$\pi^*_\kappa(\eta)=\sum_{\rho\in S_n(\lambda):\, \phi(\rho)=\eta}\pi^*_\lambda(\rho).$$
\end{cor}
\fi
%%%%%%%%%%%%%%%%%%%%%%%%%%%%%%%%%%%%%%
%%%%%%%%%%%%%%%%%%%%%%%%%%%%%%

%%%%%%%%%%%%%%%%%%%%%%%%%%%%%%%%%%%%%%%%%%%%%
%%%%%%%%%%%%%%%%%%%%%%%%%%%%%%%%%%%%%%%%%%%%%%%%%5
\iffalse 

\section{Interpolation polynomials}\label{sec:interpolation_polynomials}

\subsection{Some properties of interpolation polynomials}
\begin{definition}\label{def:int_elem}
Fix $n\geq 1$. For any $k\geq 1$, we define the \emph{elementary interpolation polynomials}
$$e^*_{k}(x_1,\dots,x_n;t):=\sum_{\begin{subarray}{c}S\subseteq \llbracket n \rrbracket\\  \#S=k\end{subarray}}\prod_{i\in S}\left(x_i-\frac{t^{\#S^c\cap \llbracket i-1\rrbracket}}{t^{n-1}}\right),$$
where the sum runs over subsets $S$ of $\llbracket n\rrbracket:=\{1,\dots, n\}$ of size $k$, and where $S^c$ denotes the complement of $S$ in $\llbracket n\rrbracket$. By convention, we set
$$e_0^*(x_1,\dots,x_n;t)=1.$$
We extend this definition to any partition $\lambda=(\lambda_1,\dots,\lambda_n)$:
$$e_\lambda^*(\bfx;t):=\prod_{i=1}^n e^*_{\lambda_i}(\bfx;t).$$
\end{definition}
The top homogeneous part of $e_{\lambda}^*$ correspond to the elementary symmetric polynomials.
We recall the following results for the interpolation ASEP polynomials; see \cite[Lemma 7.2]{BenDaliWilliams2025}.
\begin{lem}\label{lem:P_column}
        For any $\mu\in\{0,1\}^n$, we have
        \begin{equation}\label{eq:factorization_one_row}
          F_\mu^*(x_1,\dots,x_n;q,t)=\prod_{i\in S_\mu}\left(x_i-\frac{t^{\#S^c_\mu\cap \llbracket i-1\rrbracket}}{t^{n-1}}\right), \end{equation}
          where $S_\mu:=\{i\in \llbracket n\rrbracket: \, \mu_i=1\}.$
        As a consequence, 
    \begin{equation}\label{eq:P_column}
          P^*_{(1^m,0^{n-m})}(x_1,\dots,x_n;q,t)=\sum_{\mu\in S_n(1^m,0^{n-m})}F^*_{\mu}(x_1,\dots,x_n;q,t)=e^*_m(x_1,\dots,x_n;t).  
        \end{equation}
        The top homogeneous part of this equation gives
    \begin{equation}\label{eq:P_column_hom}
          P_{(1^m,0^{n-m})}(x_1,\dots,x_n;q,t)=\sum_{\mu\in S_n(1^m,0^{n-m})}F_{\mu}(x_1,\dots,x_n;q,t)=e_m(x_1,\dots,x_n;t).  
        \end{equation}
\end{lem}



\fi
%%%%%%%%%%%%%%%%%%%%%%%%%%%%%%%%%%%%%%%
%%%%%%%%%%%%%%%%%%%%%%%%%%%%%%%





%\subsection{The weak reordering property}
Fix a weakly order-preserving function $\phi:\NN\rightarrow\NN$.  Fix two partitions $\lambda$ and $\kappa$ such that $\phi(\lambda)=\kappa$.
For $\eta\in S_n(\kappa)$, define 
$$G^*_\eta(\bfx;t):=\sum_{\rho\in S_n(\lambda):\, \phi(\rho)=\eta} F^*_\rho(\bfx;1,t).$$ Let $G_{\eta}$ be the top homogeneous part of $G^*_\eta$.


The following is an analogue of \cite[Theorem 4.18]{AyyerMartinWilliams2025}, and 
can be proved in essentially the same way, using interpolation analogues of results from \cite{AlexanderssonSawhney2019}.
%on {nonsymmetric  Macdonald polynomials}. 
%\LW{!!!!!! If want to add more detail at all, then include B.3 and its proof; the others are analogues.}

\begin{thm}\label{thm:weak-reordering}
    Fix $\lambda$ and $\kappa$ as above. For all $\eta\in S_n(\kappa)$, we have at $q=1$ that 
    $$\frac{G^*_\eta(\bfx;t)}{P^*_\lambda(\bfx;1,t)}=\frac{F^*_\eta(\bfx;1,t)}{P^*_\kappa(\bfx;1,t)}.$$
\end{thm}


Given a composition $\rho$, let $\rho^-:=(\rho^-_1,\dots,\rho^-_n)$,
where
$\rho^-_i=\max(\rho_i-1,0)$.
\begin{cor}\label{cor:eta_minus}
    Consider a composition $\rho$ with $\rho_i\neq 1$ for any $1\leq i\leq n$. Let $k$ be the number of non-zero parts of $\rho$. Set $\eta=\rho^-$.
    We then have at $q=1$,
$$F^*_{\rho}(\bfx;1,t)=F^*_{\eta}(\bfx;1,t)\cdot e^*_{k}(\bfx;t).$$
\end{cor}
\begin{proof}
Let $\lambda$ and $\kappa$ be the two partitions obtained by reordering $\rho$ and $\eta$, respectively.
  Consider the weakly order-preserving function $\phi:i\mapsto \max(i-1,0)$. We then have $\phi(\rho)=\eta$. Since $\lambda$ does not have parts of size 1, and $\phi$ is bijective from $\{0,2,3,\dots\}$ to $\{0,1,2,\dots\}$, then $\rho$ is the unique composition in $S_n(\lambda)$ such that $\phi(\rho)=\eta$ and we have $G^*_\eta=F^*_\rho$. It follows then from \cref{thm:weak-reordering} that
  $$\frac{F^*_{\rho}(\bfx;1,t)}{P^*_{\lambda}(\bfx;1,t)}=\frac{F^*_{\eta}(\bfx;1,t)}{P^*_{\kappa}(\bfx;1,t)}.$$
We now recall that at $q=1$, we have from \cite{Dolega2017,BenDaliWilliams2025} that 
%For any partition $\lambda$
    \begin{equation}\label{eq:factorization}
      P^*_\lambda(x_1,\dots,x_n;1,t)=\prod_{1\leq i\leq \lambda_1}P^*_{\lambda'_i}(x_1,\dots,x_n;1,t)=\prod_{1\leq i\leq \lambda_1}e^*_{\lambda'_i}(x_1,\dots,x_n;t),
    \end{equation}
    where $\lambda'$ is the partition conjugate to $\lambda$.  Using this plus the fact that $\kappa$ is obtained from $\lambda$ by removing the largest column (of size $k$), we get that
$$\frac{P^*_{\lambda}(\bfx;1,t)}{P^*_{\kappa}(\bfx;1,t)}=e^*_{k}(\bfx;t),$$
which implies that $F^*_{\rho}(\bfx;1,t)=F^*_{\eta}(\bfx;1,t)\cdot e^*_{k}(\bfx;t).$
\end{proof}

\begin{comment}
\section{The weak reordering property for the polynomials \texorpdfstring{$h$}{h}}
We recall the following definition introduced in \cite{BenDaliWilliams2025}.
\begin{definition}\cite[Definition 4.6]{BenDaliWilliams2025}
    Given $\alpha\in \ZZ^n$, we define
    \begin{equation}\label{eq:deF_h}
      h_\alpha:=\sum_{\lambda\in \NN^n} \wt_{\alpha} a^\lambda_{\|\alpha\|}\hatfstar{\lambda^-}.
    \end{equation}    
\end{definition}
The polynomial $h_\alpha$ is the generating function for for signed multiline queues whose bottom row is Row $1'$, and that row has type $\alpha$
% \LW{rephrase?  for signed multiline queues whose bottom row is Row $1'$, and that row has type $\alpha$}.
\LW{We should make sure that we give the combinatorial interpretation of $h_{\alpha}$  in our other paper, so that we can refer to it here.}\HBD{Right! I think that these function $h_\alpha$ at $q=1$, give the stationary distribution et odd steps, only when we have exactly one vacancy (I will have to double check that). This is related to the fact that Corollary 2.2 is only true under the condition that there is a unique part of size $0$ in $\mu$. If we decide to include this, we will need the combinatorial interpretation of the polynomials $h_\alpha$.}

In the following, we assume that $\lambda$ and $\mu$ have a unique part of size 0.
For example, $\lambda=(3,2,2,1,1,0)$ and $\phi(0)=0$, $\phi(2)=\phi(3)=2$ and $\phi(1)=1$.
We extend the function $\phi:\NN\rightarrow \NN$, to an odd function $\phi:\ZZ\rightarrow\ZZ$ ($\phi(n)=-\phi(n)$).
\begin{cor}
     For any signed permutation $\alpha\in S_n^{\pm}(\mu)$, we have at $q=1$
    $$\frac{\sum_{\beta\in S_n^{\pm}(\lambda):\phi(\beta)=\alpha}h_\beta}{P^*_\lambda}=\frac{h^*_\alpha}{P^*_\mu}.$$
\end{cor}
\textit{Sketch of the proof:}
    We start with the packed case (\ie when $\lambda_n=0$). In this case, the polynomials $h_\alpha$ are directly obtained from $F^*_\eta$ by
    \begin{equation*}
        h_\alpha=\wt_{\alpha}\frac{F^*_{\mu}}{\prod_{1\leq i\leq n-1}(x_i-\frac{1}{t^{n-1}})}
    \end{equation*}
    and the corollary follows from \cref{thm:weak-reordering}. We obtain the general case by application of the Hecke operators.


% The action of the Hecke operators on $h_\alpha$ was given in \cite[Lemma 4.7]{BenDaliWilliams2025}. The following particular case will be useful.
% \begin{lem}
    
% \end{lem}
    
\end{comment}

%\section{The proof of the main theorem}\label{sec:proof}
%\subsection{The proof of \texorpdfstring{\cref{thm:main}}{the main theorem} for restricted partitions}
%when \texorpdfstring{$\lambda$}{lambda} has distinct parts}

%%%%%%%%%%%%%%%%%%%%%%%%%%%%%%%%%%%
\iffalse%%%%%%%%%%%%%%%%%%%%
\begin{definition}\label{def:restricted}
    A \emph{restricted composition}
    (respectively, a \emph{restricted partition}) $\mu=(\mu_1,\dots,\mu_n)$ is a composition (respectively, partition)
    which has all parts distinct, a unique part of size $0$, and no part of size $1$.
\end{definition}

\emph{In this section we will assume that each $\lambda$ is a restricted partition as in \cref{def:restricted}.}
\fi
%%%%%%%%%%%%%%%%%%%%%%%%%%%%%55


%%%%%%%%%%%%%%%%%%%%%%%%%%%%%%%%%%%%%%%%%%%%
%%%%%%%%%%%%%%%%%%%%%%%%%%%%%%%%%%%%%%%%%%
\iffalse
The following lemma relates the transition probabilities of the classical $t$-Push TASEP (or equivalently \ref{Step1} in \cref{def:intpushTASEP}) to the generating function of two-line queues.


\begin{prop}[{\cite[Lemma 5.4]{AyyerMartinWilliams2025}}]\label{prop:Pone}
Fix $\mu,\rho\in S_n(\lambda)$, and let $j$ be the index such that $\rho_j=0$.  Then the following quantities are equal:
\begin{itemize}
\item the probability that in 
Step 1 of \cref{def:intpushTASEP} (i.e. the $t$-Push TASEP) we transition from $\mu$ to $\rho$ when the bell rings at site $j$ in $\mu$,
\item the weight generating function of classical
two-line queues with top row $\mu$ and bottom row $\rho$.
\end{itemize}
That is,
$$\PPone{_j}(\mu,\rho)=a^\mu_{\rho}.$$
\end{prop}


To help prepare the reader for the proof of \cref{prop:Ptwo}, we 
will prove \cref{prop:Pone} here.  The main idea  is that
a two-line queue encodes a transition of the $t$-Push TASEP.  For instance, \cref{fig:two_line_queue_restricted} encodes 
Step 1 in the transition illustrated in \cref{fig:Markov_chain_2}:
the top row encodes the initial state, each nontrivial pairing encodes a particle which moves, and 
the bottom row encodes the state after Step 1.

\begin{figure}[t]
    \centering
    \includegraphics[width=0.5\linewidth]{2_line_queue_restricted.pdf}
    \caption{The (classical) two-line queue encoding Step 1 in the transition of \cref{fig:Markov_chain_2}.}
    \label{fig:two_line_queue_restricted}
\end{figure}




\begin{proof}
Note that since by assumption $\mu$ and $\rho$ are restricted compositions, they uniquely determine the pairings of the balls in a two-line queue whose top row is $\mu$ and bottom row is $\rho$.  Thus, there is at most one two-line queue with top row $\mu$ and bottom row $\rho$.  There exists one such queue if and only if when we pair the balls with the same labels, %in a potential two-line queue with top row $\mu$ and bottom row $\rho$, 
the bottom ball in a nontrivial pairing always has 
a ball with a smaller label directly above it.
 This condition is equivalent to the statement (*) that if we write the composition $\rho$ underneath the composition $\mu$, then for each position $i\neq j$ such that $\mu_i \neq \rho_i$, we have $\mu_i < \rho_i$.
%connecting
%balls labeled $b$, the ball directly above the bottom ball labeled $b$ has a label $a<b$.
Moreover, when we have such a two-line queue,
its weight is the product of the pairing weights, each of which have the form 
$$\frac{(1-t)t^{\skipped(p)}}{1-t^{\free(p)}} = \frac{t^{\skipped(p)}}{[\free(p)]_t}.$$

Meanwhile, there is a unique series of ball displacements 
that must occur in order to realize a transition
in the $t$-Push TASEP from $\mu$ to $\rho$: namely, 
if $b>b^{(1)} > \dots > b^{(\ell)}$ are the letters in 
$\mu$ which appear in different locations in $\rho$, 
where $b$ is the ball at location $j$ in $\mu$, then 
first the ball $b$ must travel to the location of ball
$b^{(1)}$ (settling in that position and displacing 
the ball labeled $b^{(1)}$); then the ball $b^{(1)}$
must travel to the location of ball $b^{(2)}$ (settling in that position and displacing the ball labeled $b^{(2)}$); and so on.  If the resulting configuration is $\rho$, then there is a nonzero probability of transitioning from $\mu$ to $\rho$. 
Note e.g. that if $j_1$ is the location of ball $b^{(1)}$ in $\mu$,
then we are saying that $\mu_{j_1} = b^{(1)}$ and $\rho_{j_1}=b > \mu_{j_1}$ (and similarly for subsequent balls), which is equivalent to the statement (*) above.  Thus, we can construct a two-line queue whose pairings encode the sequence of ball displacements.
Moreover, note that in Step 1 of \cref{def:intpushTASEP},
the probability of ball $b$ displacing ball $b^{(1)}$ 
is $\frac{t^{k-1}}{[m]_t}$, when there are $m$ weaker particles and $b^{(1)}$ is the $k$th such particle that $b$ encounters when it travels clockwise; this exactly corresponds to the statistic 
$\frac{(1-t)t^{\skipped(p)}}{1-t^{\free(p)}} = \frac{t^{\skipped(p)}}{[\free(p)]_t}$ 
in the corresponding two-line queue.
\end{proof}

Recall that \cref{fig:two_line_queue_restricted} shows the two-line queue corresponding to Step 1 in the transition of \cref{fig:Markov_chain_2}, where $j=5$, $\mu=(3,0,6,8,7,4,5,2)$ and $\rho=(5,3,6,8,0,4,7,2)$. One then checks that 
$$\PPone{j}(\mu,\rho)=\frac{t}{[6]_t}\times\frac{t}{[4]_t}\times \frac{1}{[2]_t}=a^\mu_\rho.$$



\begin{remark}\label{rmk:subsets}
Each possible transition of the t-Push TASEP and each instance of Step 2 above can be indexed by 
choosing a subset of positions, where the corresponding particles have distinct labels.  For the $t$-Push TASEP, once we have the subset of positions, we know that the particles have to get bumped in decreasing order of their labels.  And in Step 2, the particles get bumped from left to right.
\end{remark}



\begin{comment}
In the following, we assume that $\lambda$ is a partition with distinct parts and one part of size 0.
\begin{notation}\label{not:m_set}
    Fix $\kappa\in S_n(\lambda)$ and $j$ (the unique index) such that $\kappa_j=0$.
    \begin{itemize}
        \item Let  $\MM\subseteq\llbracket n\rrbracket$ be  a subset satisfying $j=\max(\MM)$. We associate to $T$ the composition $\nu$ $\kappa$ by moving the particles indexed by $\MM$ as explained in \cref{rmk:subsets}, i.e
    \begin{equation*}
    \begin{cases}
        \nu_i=\kappa_i& \text{if $i\notin\MM$,}\\
         \nu_{i}=\kappa_j=0 &\text{if $i=\min(\MM)$,}\\
         \nu_{i}=\kappa_{\max(\MM\cap \llbracket i-1\rrbracket)}& \text{if $i\in \MM$ and $i\neq \min(\MM)$.}
    \end{cases}
    \end{equation*}
     \item Conversely, if $\nu$ is a composition that can be obtained from $\kappa$ by \ref{Step2} (i.e $\PPtwo{j}(\kappa,\nu)\neq 0$) then associate to it the set 
    $$\MM:=\{i\in \llbracket n\rrbracket:\ \kappa_i\neq \nu_i\}\cup \{j\}.$$
    \end{itemize}
   Under the assumption that $\kappa$ has distinct parts, the two maps described above are inverse of each other.
\end{notation}
\begin{example}
    If $n=7$, $\kappa=(1,4,6,5,0,2,3)$, $j=5$ and $\MM=\{1,2,4,5\}$, then the composition $\nu$ given by \cref{not:m_set} is $\nu=(0,1,3,4,5,2,3).$
\end{example}


We  have the following explicit expression for the transition probabilities $\PPtwo{}$ which follows from the definition.
\begin{lem}
Let $\kappa$, $j$, $T$ and $\nu$ as in \cref{not:m_set}.
We then have
    $$\PPtwo{j}(\kappa,\nu)=\prod_{k<j:\, k\in \MM} \mfr_k\prod_{k<j:\, k\notin \MM}(1-\mfr_k),$$
where 
$$\mfr_k=
\begin{cases}
    \mfp_k &\text{if $\kappa_k> \nu_k$,}\\
    \mfq_k &\text{if $\kappa_k<\nu_k$}.
\end{cases}$$
\end{lem}    
\end{comment}

% We now write the transition rates using the generating functions of signed two-line queues.

We now give an analogous result relating the transition probabilities of \ref{Step2} of \cref{def:intpushTASEP}  to the generating function for signed two-line queues.
\fi
%%%%%%%%%%%%%%%%%%%%%%%%%%%%%%%%%%%%%%%5
%%%%%%%%%%%%%%%%%%%%%%%%%



\begin{prop}\label{prop:Ptwo}
     Fix $\rho,\nu\in S_n(\lambda)$, and let $j$ be the index such that $\rho_j=0$. We have 
     $$\PPtwo{j}(\rho,\nu)=\frac{c^\rho_\nu}{\prod_{k<j}\left(x_k-\frac{1}{t^{n-2}}\right)\prod_{k>j}\left(x_k-\frac{1}{t^{n-1}}\right)},$$
     or equivalently,
     $$ P_j\cdotp \PPtwo{j}(\rho,\nu)=\frac{c^\rho_\nu}{e_{n-1}^*},$$
     where $c^\rho_\nu$ is the coefficient from \cref{eq:def_c}, % and \cref{lem:c}, 
     i.e. 
     %the generating function for signed
     %two-line queues. \HBD{
     the generating function for the set $\Gbar_{\nu}^\rho$.
\end{prop}


%%%%%%%%%%%%%%%%%%%%%%%%%%
\iffalse 
\begin{figure}
    \centering
    \includegraphics[width=0.5\linewidth]{unsigned_line_queue_restricted.pdf}
    \caption{The element of $\Gbar^{(5,3,6,8,0,4,7,2)}_{(0,3,6,5,8,4,7,2)}$ (see \cref{def:Gbar}) 
    %paired ball system \LW{do we want to talk here about the paired ball system or the signed two-line queue?  for instance, it is important here that we don't allow wrapping around} 
    encoding the Step 2 of the transition of \cref{fig:Markov_chain_2}.}
    \label{fig:unsigned_two_line_queue_restricted}
\end{figure}
\fi%%%%%%%%%%%%%%%%%%%%%
%%%%%%%%%%%%%%%%%%%%%%%%%%%%%%%%


%As in the proof  of \cref{prop:Pone}, 
The  idea of the proof below is that a signed two-line queue encodes Step 2 of the 
interpolation $t$-Push TASEP.  
%For instance, 
%\cref{fig:unsigned_two_line_queue_restricted} encodes Step 2 
%of the transition of \cref{fig:Markov_chain_2}.

\begin{proof}
 Note that \ref{Step2} of \cref{def:intpushTASEP} is encoded by 
an element of a set $\Gbar^{\rho}_{\nu}$ (see \cref{def:Gbar}).  
% the paired ball systems \LW{call it paired ball system or signed two-line queue?} described in \cref{ssec:MLQ_distinct_parts}. 
Indeed, the transition in \ref{Step2} from the configuration $\rho$ to the configuration $\nu$ is possible if and only there is an
element $\Qbar$ in $\Gbar^{\rho}_\nu$ (recall that this set contains at most one element).
More precisely,  a particle labeled $a>0$ which moved from position $k\in \llbracket n\rrbracket$ to a position $k'$, corresponds to a non trivial pairing in $\Qbar$ connecting a ball labeled $a$ in column $k$ of the top row to a ball labeled $a$ in column $k'$ of the bottom row. Particles which do not move correspond to trivial pairings. 

We now claim that $\wt(\Qbar)$ divided by 
$D:=\prod_{k<j}\left(x_k-\frac{1}{t^{n-2}}\right)\prod_{k>j}\left(x_k-\frac{1}{t^{n-1}}\right)$ gives $\PPtwo{j}(\rho,\nu)$.  We will prove the claim below by showing that each 
ball or pairing weight in $\wt(\Qbar)$,
divided by one of the factors in $D$, equals one of the skipping/ displacement
probabilities from \cref{Step2}
%\eqref{eq:probskip} or \eqref{eq:prob} 
(whose product is $\PPtwo{j}(\rho,\nu)$).  Note that in what follows,
instead of associating the weight 
$(1-t)t^{\skipped(p)}$ to each nontrivial pairing,
we will associate $(1-t)$ to the top ball in each nontrivial pairing,
and a factor of $t$ to each skipped ball.
\begin{itemize}
    \item Each ball in column $k>j$ of $\Qbar$ is necessarily trivially paired, since no ball in position $k>j$ get skipped or displaced in \ref{Step2}.
      In $\Qbar$ this ball gets  weight $x_k-\frac{1}{t^{n-1}}$; when we divide this weight by the $k$th factor of $D$, we get $1$, which corresponds to the fact that balls in position $k>j$ do not contribute to $\PPtwo{j}(\rho,\nu)$.
        \item A ball in $\Qbar$ labeled $b$ in column $k<j$ which is trivially paired, and which is not skipped by a ball $a>b$, also has weight $x_k-\frac{1}{t^{n-1}}$. When we divide this weight by the $k$th factor of $D$, we get $1-\mfp_k$ (see \eqref{eq:1minuspk}).  This is what we desired, because such a trivial pairing in $\Qbar$ corresponds to a particle labeled $b$ which is skipped over by a particle with a smaller label, and hence contributes $1-\mfp_k$
        to $\PPtwo{j}(\rho,\nu)$.
    \item A ball in $\Qbar$ labeled $b$ in column $k<j$ which is trivially paired, and which is skipped by a ball $a>b$, gets a weight $t(x_k-\frac{1}{t^{n-1}})$.  When we divide this weight by the 
    $k$th factor of $D$, we get $1-\mfq_k$
    (see \eqref{eq:1minuspk}).
    This is what we desired, because such a trivial pairing corresponds to a particle labeled $b$ skipped over by a particle with a larger label, and hence contributes $1-\mfq_k$ to $\PPtwo{j}(\rho,\nu)$.
    \item A ball labeled $b$ in the top row of $\Qbar$ in column $k<j$ which has a ball labeled $a<b$ below it gets a weight $(1-t)\frac{1}{t^{n-1}}$
    % \LW{This should be $(1-t)\frac{1}{t^{n-1}}$, right?} 
    (the factor $(1-t)$ is the nontrivial pairing weight). When we divide this weight by the $k$th factor of $D$, we get $\mfp_k$.  This is what we desired, because this pairing corresponds to a particle labeled $b$ being displaced by a particle with a smaller label, and hence contributing $\mfp_k$ 
    to $\PPtwo{j}(\rho,\nu)$.
    \item A ball labeled $b$ in the top row of 
    $\Qbar$ in column $k<j$ which has a ball labeled $a>b$ below it gets a weight $(1-t)x_k$ (the factor $(1-t)$ is the nontrivial pairing weight). 
     When we divide this weight by the $k$th factor of $D$, we get $\mfq_k$.  This is what we desired, because this pairing
    corresponds to a particle labeled $b$ being displaced by a particle with a larger label, and hence contributing $\mfq_k$ to $\PPtwo{j}(\rho,\nu)$.\qedhere
\end{itemize}
\end{proof}

%%%%%%%%%%%%%%%%%%%%%%%%%%%%%%%%%%%%%%%%%%%%
%%%%%%%%%%%%%%%%%%%%%%%%%%%%%%%%%%%%%%%
\iffalse 

Recall that \cref{fig:unsigned_two_line_queue_restricted} shows the paired ball system corresponding to Step 2 in the transition of \cref{fig:Markov_chain_2}, where $j=5$, $\rho=(5,3,6,8,0,4,7,2)$ and $\nu=(0,3,6,5,8,4,7,2)$. We then have 
\begin{align*}
  \PPtwo{5}(\rho,\nu)&=\mfp_1\cdotp(1-\mfq_2)\cdotp(1-\mfp_3)\mfp_4\\
  &=\frac{t^{-7}(1-t)}{x_1-t^{-6}}\cdotp \frac{t (x_2-t^{-7})}{x_2-t^{-6}}\cdotp\frac{x_3-t^{-7}}{x_3-t^{-6}}\cdotp\frac{t^{-7}(1-t)}{x_4-t^{-6}}.
\end{align*}
On the other hand, for the paired ball system $\Qbar$ of \cref{fig:unsigned_two_line_queue_restricted}, has ball weight (as given in \cref{eq:weights_unsigned_balls}) 
$$t^{-7}\cdotp (x_2-t^{-7})\cdotp(x_3-t^{-7})\cdotp t^{-7}\cdotp (x_6-t^{-7})\cdotp (x_7-t^{-7}) \cdotp(x_8-t^{-7})$$
and its pairing weight (defined in \cref{eq:pairing_weight_unsigned}) is 
$$(1-t)^2\cdotp t,$$
and $\wt(\Qbar)$ is the product of these two quantities. We then get that
$$\frac{c^\rho_\nu}{\prod_{k<5}\left(x_k-t^{-6}\right)\prod_{k>5}\left(x_k-t^{-7}\right)}
=\frac{\wt(\Qbar)}{\prod_{k<5}\left(x_k-t^{-6}\right)\prod_{k>5}\left(x_k-t^{-7}\right)}=\PPtwo{5}(\rho,\nu).$$
\fi
%%%%%%%%%%%%%%%%%%%%%%%%%%%%%%%%%%%%%
%%%%%%%%%%%%%%%%%%%%%%%%%%%%%%%%%%
%We then obtain the following result.
%get the following proposition connecting the interpolation $t$-Push TASEP model to signed multiline queues.
\begin{prop}\label{prop:intTASEP_SMLQ}
If $\lambda$ is restricted, and $\mu,\nu\in S_n(\lambda)$, then    
\begin{align*}
    \PP(\mu,\nu)=\sum_{\rho\in S_n(\lambda)}  \frac{a^\mu_\rho c^\rho_{\nu}}{e^*_{n-1}}.
\end{align*}
\end{prop}
\begin{proof}
    Combining 
    \cite[Lemma 5.4]{AyyerMartinWilliams2025}
     and \cref{prop:Ptwo}, we get
    \begin{align*}
    \PP(\mu,\nu)
    &=\sum_{1\leq j\leq n}P_j\sum_{\rho\in S_n(\lambda):\rho_j=0}\PPone{j}(\mu,\rho)\PPtwo{j}(\rho,\nu)\\
    &=\sum_{1\leq j\leq n}\sum_{\rho\in S_n(\lambda):\rho_j=0}  \frac{a^\mu_\rho c^\rho_{\nu}}{e^*_{n-1}}\\
    &=\sum_{\rho\in S_n(\lambda)}  \frac{a^\mu_\rho c^\rho_{\nu}}{e^*_{n-1}}.\qedhere
\end{align*}
\end{proof}

%\subsubsection{The proof of \cref{thm:main}}
\begin{proof}[Proof of \cref{thm:main}]
Fix a restricted partition $\lambda$.
%with distinct parts, with one part of size 0 and no parts of size 1. 
%\HBD{Should we say restricted here?}
Let $\nu\in S_n(\lambda)$.
From 
\cite[Theorem 1.15 and Lemma 5.6]{BenDaliWilliams2025}, 
we have
 \begin{equation*}
   F^*_\nu(\bfx;1,t)=\sum_{\eta\in\NN^n}{F^{*\eta}_\nu}(\bfx;t) F^*_{\eta^{-}}(\bfx;1,t), 
 \end{equation*}
 where 
 $$F^{*\eta}_{\nu}(\bfx;t) :=\sum_{\alpha\in\ZZ^n}b_\nu^\alpha \wt_\alpha a^\eta_{\lVert \alpha\rVert}=\sum_{\kappa\in \NN^n}a_\kappa^\eta c^\kappa_\nu .$$
But we know from 
% \LW{Oops!  We're relying on something that comes later}
\cref{cor:eta_minus} that
$$F^*_{\eta^-}(\bfx;1,t)=\frac{F^*_{\eta}(\bfx;1,t)}{e_{n-1}^*(\bfx;t)},$$
we use here the fact that $\eta$ has a unique part of size 0.
% \HBD{Change $q=1$ by $(\bfx,;1,t)$.}
Hence
\begin{equation*}
   F^*_\nu(\bfx;1,t)=\sum_{\eta\in\NN^n}F^*_{\eta}(\bfx;1,t)\sum_{\kappa\in \NN^n}\frac{a_\kappa^\eta c^\kappa_\nu}{e_{n-1}^*(\bfx;t)},
 \end{equation*}
which can be rewritten using the transition probabilities of the interpolation $t$-Push TASEP (\cref{prop:intTASEP_SMLQ}) we get
\begin{equation*}
   F^*_\nu(\bfx;1,t)=\sum_{\eta\in\NN^n}F^*_{\eta}(\bfx;1,t)\PP(\eta,\nu).
 \end{equation*}
This proves that $F^*_\mu(\bfx;1,t)$ are proportional to the stationary distribution of the interpolation $t$-Push TASEP $\pi^*_\lambda(\mu)$.
Finally, we use 
the fact that 
$P^*_\lambda=\sum_{\mu\in S_n(\lambda)}F^*_\mu$ to deduce that $\frac{F^*_\mu(\bfx;1,t)}{P^*_\lambda(\bfx;1,t)}=\pi^*_\lambda(\mu).$
% \HBD{Or maybe we should use $\pi^*_\lambda(\mu)$ to denote the stationary distribution of the interpolation $t$-Push TASEP?} \LW{Yes, I like that notation!}
\end{proof}



%%%%%%%%%%%%%%%%%%%%%%%%%%%%%%%%%%%%%%%%%%%%%%%
%%%%%%%%%%%%%%%%%%%%%%%%%%%%%%%%%%%%%%%%%%%%%%%%
\iffalse 
\subsection{Proof of \texorpdfstring{\cref{thm:main}}{the main theorem} for arbitrary partitions}\label{sec:arbitrary}
\begin{proof}[Proof of \cref{thm:main}]
Fix a partition $\kappa=(\kappa_1,\dots,\kappa_n)$, and $\eta\in S_n(\kappa)$.  We want to prove that 
$$\pi_\kappa^*(\eta)=\frac{F^*_\eta(\bfx;1,t)}{P^*_\kappa(\bfx;1,t)}.$$
It is not hard to see that we can find 
a restricted partition $\lambda$ and  
%be a restricted partition, and let 
a weakly order-preserving function $\phi$ such that $\phi(\lambda)=\kappa$.
By \cref{cor:lumping}, 
%From the lumping property of the  interpolation %\texorpdfstring{$t$}{t}-Push TASEP \cref{prop:lumping}, 
we know that
        $$\pi^*_\kappa(\eta)=\sum_{\rho\in S_n(\lambda):\, \phi(\rho)=\eta}\pi^*_{\lambda}(\rho).$$
Since \cref{thm:main} holds for any restricted partition $\lambda$,
%We apply the result of the previous section on $\lambda$, 
we get
$$\pi^*_\kappa(\eta)=\sum_{\rho\in S_n(\lambda):\, \phi(\rho)=\eta}\frac{F^*_\rho(\bfx;1,t)}{P^*_\lambda(\bfx;1,t)}.$$
Now from \cref{thm:weak-reordering}, we have
$$\sum_{\rho\in S_n(\lambda):\, \phi(\rho)=\eta}\frac{F^*_{\rho}(\bfx;1,t)}{P^*_\lambda(\bfx;1,t)}
=\frac{G^*_{\eta}(\bfx;t)}{P^*_\lambda(\bfx;1,t)}
=\frac{F^*_\eta(\bfx;1,t)}{P^*_\kappa(\bfx;1,t)}.$$
We then deduce the desired statement
\begin{equation*}
  \pi^*_\kappa(\eta)=\frac{F^*_\eta(\bfx;1,t)}{P^*_\kappa(\bfx;1,t)}.\qedhere  
\end{equation*}
\end{proof}

\section{Density formulas}\label{sec:density}

The following proposition is a consequence of
\cref{thm:main} and 
\cite[Lemma 7.2]{BenDaliWilliams2025}\HBD{and \cref{lem:P_column}}.
\begin{prop}\label{prop:01}
Let $\lambda = (1^{m_1}, 0^{m_0})$ be a partition with $n$ parts whose parts are all $1$'s or $0$'s.  The stationary probability $\pi_\lambda^*(\mu)$ of $\mu \in S_n(\lambda)$ for the 
interpolation $t$-Push TASEP is 
$$\pi^*_\lambda(\mu) = 
\frac{F_{\mu}^*(x_1,\dots,x_n; q, t)}{e^*_{m_1}(x_1,\dots,x_n;t)} =
\frac{1}{e^*_{m_1}(x_1,\dots,x_n;t)} \prod_{i\in S_\mu}\left(x_i-\frac{t^{\#S^c_\mu\cap \llbracket i-1\rrbracket}}{t^{n-1}}\right), $$
    where $S_{\mu}:=\{i: \mu_i=1\}.$
\end{prop}
Recall that the \emph{density} at site $i$ is the probability of finding a particle at site $i$ in the stationary distribution.
We  use angle brackets 
$\langle \cdot \rangle$ to denote expectations in the stationary distribution.  Let 
$\eta_j$ be the indicator variable for a particle at site $j$, so that 
$\langle \eta_j\rangle$ is the density at site $j$.  
%\LW{In Ayyer-Martin Corollary 5, they immediately get a formula for density at site $j$, because of the factorization formula (top homogeneous part of factorization above).  However, in most cases we do not get a straightforward analogue of their formula because it is not the case that all $F^*_{\mu}$ where $\mu_j=1$ share the same factor $(x_j-something)$.  But if $j=1$ or $n$, the factor is the same.}

From the definition of density, we obtain the following formula. 
\begin{lem}\label{lem:density}
Let $\lambda = (1^{m_1}, 0^{m_0})$ be a partition with $n$ parts whose parts are all $1$'s or $0$'s.
  The density at state $j$ in the stationary distribution is given by 
  $$\langle \eta_j \rangle =
  \frac{\sum F^*_{\mu}(x_1,\dots,x_n; q,t)}{e^*_{m_1}(x_1,\dots,x_n;t)},$$
  where the sum is over all $\mu\in S_n(\lambda)$ which have a $1$ in site $j$.
\end{lem}
Using \cref{prop:01} and \cref{lem:density}, we obtain the following.
\begin{cor}\label{cor:density}
    Let $\lambda = (1^{m_1}, 0^{m_0})$ be a partition with $n$ parts whose parts are all $1$'s or $0$'s.
  The density at state $j=1$ in the stationary distribution is given by 
  \begin{equation}\label{eq:density}
  \langle \eta_1 \rangle =
  \frac{(x_1-\frac{1}{t^{n-1}}) e^*_{m_1-1}(tx_2, tx_3,\dots,tx_n;t)} {t^{m_1-1} e^*_{m_1}(x_1,\dots,x_n;t)},
  \end{equation}
   and the density at state $j=n$ in the stationary distribution is given by 
  $$\langle \eta_n \rangle =
  \frac{(x_n-\frac{t^{m_0}}{t^{n-1}}) e^*_{m_1-1}(tx_1, tx_2, \dots,tx_{n-1};t)} {t^{m_1-1} e^*_{m_1}(x_1,\dots,x_n;t)},$$
\end{cor}
% \HBD{In general, can we write 
% \begin{align*}
%   \langle \eta_j \rangle&=\frac{1}{e^*_{m_1}(x_1,\dots,x_n)}\sum_{0\leq k\leq \min(m_1,j)}\\
%   &\frac{1}{t^{k+n-j-1}}\left(x_j-\frac{t^{j-k-1}}{t^{r}}\right)e_{k}(t^{n-j+1}x_1,\dots, t^{n-j+1}x_{j-1})e_{m_1-k-1}(t^{k}x_{j+1},\dots, t^{k}x_{n}).
% \end{align*}
% with $r=(j-1)(n-j+1)+k(m_1-k-1)$. Maybe to avoid this kind of extra powers of t, we can introduce for this section a slightly more general notation
% $$e_{k,r}^*(x_1,\dots,x_n):=\sum_{\begin{subarray}{c}S\subseteq \llbracket n \rrbracket\\  \#S=k\end{subarray}}\prod_{i\in S}\left(x_i-\frac{t^{\#S^c\cap \llbracket i-1\rrbracket}}{t^{r}}\right).$$
% for any $n,r\geq 0,$ and $0\leq \ell\leq n$.
% }
%\LW{Do you want to keep the above or not?  I don't have a strong opinion here.} \HBD{I removed the general expressions and added the following comment:}
One can actually write an explicit expression for $\langle\eta_j\rangle$ for any $1\leq j\leq n$. Here we only give the expressions for $j=1$ and $j=n$, for which the formulas are slightly simpler.

\HBD{Added this}
We will use the notation $\bfx:=(x_1,\dots,x_n)$, $\bfx^{(1)}:=(x_2,\dots,x_n)$,  $t\bfx:=(tx_1,\dots,tx_n)$ and so on.
\begin{definition}\label{def:sstar}
For any partition $\lambda=(\lambda_1,\dots,\lambda_n)$, we define the \emph{$t$-interpolation Schur polynomial 
% \LW{the next few $s$'s were missing stars, so I added them} 
$s_\lambda^*(x_1,\dots,x_n;t)=s^*_\lambda(\bfx;t)$} by
    $$s^*_\lambda(\bfx;t)=P_\lambda(\bfx;t,t).$$
\end{definition}

\HBD{Would it be better if we move \cref{lem:2column} here, since we are using it in the the proof of the proposition?}

The following result is an analogue of \cite[Corollary 7.1]{AyyerMartinWilliams2025}.  

\begin{prop}
Recall, as in \cref{not:type}, that  $M_i = m_i+m_{i+1} + \dots + m_L.$  Let $s^*_{\lambda}$ be the $t$-interpolation Schur polynomial from \cref{def:sstar}.
   The density of the particle of species $i$ in the first site in 
   the interpolation $t$-Push TASEP with content $\lambda=(L^{m_L}, \dots, 1^{m_1}, 0^{m_0})$ is given by
   % \HBD{Should we add the parameter $t$ in the $t$-elementary and $t$-Schur functions below, or simply add that $$e_k(\bfx)\equiv e_k(\bfx;t).$$
   % Also, would $\bfx^{(1)}$ be a good notation for $(x_2,\dots,x_n)$? or maybe $\bfx_{2;n}$?}
   \small{
   \begin{align*}
   %\label{eq:density}
       \langle \eta_1^{(i)} \rangle =&
       \frac{(x_1-\frac{1}{t^{n-1}}) \left(e_{M_i-1}^*\left(t\bfx^{(1)};t\right) e^*_{M_{i+1}}\left(\bfx^{(1)};t\right)  - t^{m_i} e_{M_{i+1}-1}^*\left(t\bfx^{(1)};t\right) e_{M_i}^*\left(\bfx^{(1)};t\right) \right)}{t^{M_{i}-1} e^*_{M_i}\left(\bfx;t\right) e^*_{M_{i+1}}\left(\bfx;t\right)} \\
       =&
       \frac{(x_1-\frac{1}{t^{n-1}}) 
       \ s^*_{(2^{M_{i+1}},1^{m_i-1})}\left(t\bfx^{(1)};t\right)
       }{t^{M_{i+1}+M_i-1} e^*_{M_i}\left(\bfx;t\right)e^*_{M_{i+1}}\left(\bfx;t\right)}.
   \end{align*}}
\end{prop}

\begin{proof}
By \cref{prop:lumping} \HBD{and the lumping property (\cref{cor:lumping})},  the density of the particle
    of species $i$ is the density of the particle of species $1$ in the single species
    Markov chain with $M_i$ particles, minus the density of the particle of species $1$
    in the single species Markov chain with $M_{i+1}$ particles.
    Thus using \cref{eq:density}, we have that 
     \begin{align*}
       \langle \eta_1^{(i)} \rangle &= 
       \frac{(x_1-\frac{1}{t^{n-1}}) e^*_{M_i-1}\left(t\bfx^{(1)};t\right)}{t^{M_i-1} e^*_{M_i}(\bfx;t)}  - \frac{(x_1-\frac{1}{t^{n-1}}) e^*_{M_{i+1}-1}\left(t\bfx^{(1)};t\right)}{t^{M_{i+1}-1} e^*_{M_{i+1}}(\bfx;t)} \\
       &= \frac{(x_1-\frac{1}{t^{n-1}})}{t^{M_{i}-1}} \cdot 
       \frac{e^*_{M_i-1}\left(t\bfx^{(1)};t\right) e^*_{M_{i+1}}(\bfx;t)-t^{m_i} e^*_{M_{i+1}-1}\left(t\bfx^{(1)};t\right) e^*_{M_i}(\bfx;t)}{e^*_{M_i}(\bfx;t) e^*_{M_{i+1}}(\bfx;t)}, \label{eq:sofar}
   \end{align*}
   where $\bfx:=(x_1,\dots,x_n)$.
If we now use the fact that 
\begin{equation}
e^*_d(\bfx;t) = \left(x_1-\frac{1}{t^{n-1}}\right) t^{-(d-1)} e^*_{d-1}\left(t\bfx^{(1)};t\right) + e^*_d\left(\bfx^{(1)};t\right),    
\end{equation}
we get the first statement of the proposition.  To get the second statement, we use \cref{lem:2column}.
\end{proof}

%\LW{Dual Jacobi-Trudi says that 
%$$s_{(2^a, 1^b)} = e_{a+b} e_a - e_{a+b+1} e_{a-1}$$
%(where everything is a polynomial in $x_1,\dots,x_n$.
%Now we'd like an analogous formula on the interpolation %level that we can apply to \eqref{eq:density}: the RHS %of such a dual JT formula would be
%$$e^*_{a+b}(tx_1,\dots,tx_n) e^*_a(x_1,\dots,x_n) - %t^{b+1} e^*_{a-1}(tx_1,\dots,tx_n) e^*_{a+b+1}(x_1,%\dots,x_n).$$}





\begin{appendices}
    
\section{Shape permuting operators and Knop–Sahi recurrence}\label{app:shape_permuting}
The main purpose of \cref{app:shape_permuting} and \cref{app:permuted_basement} is to prove \cref{thm:weak-reordering}. We will follow the strategy of \cite{AlexanderssonSawhney2019} and \cite{AyyerMartinWilliams2025} (which proved the homogeneous case of \cref{thm:weak-reordering}). 

Fix $n\geq 1$. Let $\YY_n$ denote the set of integer partitions 
$\lambda=(\lambda_1,\dots,\lambda_n)=(\lambda_1 \geq \dots \geq \lambda_n)$ with at most $n$ parts.  We let $|\lambda|$ denote the sum $\lambda_1+ \dots + \lambda_n$ of the parts of the partition and call it the \emph{size} of $\lambda$.
Let $\mcP_n$ denote the ring of polynomials in $n$ variables, and 
let $\mcP_n^{(d)}$ denote the polynomials of degree at most $d$. 


\subsection{The dehomogenization operator \texorpdfstring{$\psi$}{psi}}
For any partition $\lambda\in\YY_n$ of size $d$, we define the space $\VV\subset\mcP^{(d)}_n$ by
$$\VV:=\left\{f\in \mcP_n^{(d)}: f(\widebar\nu)=0 \text{ for any $|\nu|\leq |\lambda|$ and $\nu\notin S_n(\lambda)$}\right\},$$
where $\widebar{\nu}$ is the sequence defined in \cref{eq:k2}.

We recall that the \textit{nonsymmetric Macdonald polynomials} $E_{\mu} = E_\mu(\bfx;q,t)$, where $\mu$ ranges over all compositions with at most $n$ parts, form a basis for the space of polynomials $\mcP_n$; see e.g. \cite{HaglundHaimanLoehr2008}. 
We will also use the interpolation analogue of these polynomials $E^*_\mu(\bfx;q,t)$ \cite{Knop1997b,Sahi1996}, which can be characterized as follows. 
\begin{thm}[\cite{Knop1997b,Sahi1996}]\label{thm:Knop_Sahi}
    Fix $\mu\in \NN^n$ of size $d$. There exists a unique polynomial $E^*_\mu\in\mcP_n^{(d)}$, called the \emph{nonsymmetric interpolation Macdonald polynomial}, such that
    \begin{itemize}
        \item $[x^\mu]E^*_\mu=1$ (so in particular, $E^*_{\mu}$ has degree $d$),
        \item $E^*_\mu(\widetilde\nu)=0$ for any $\nu\in \NN^n$ satisfying $|\nu|\leq d$ and $\nu\neq \mu$.
    \end{itemize}
    Moreover, the top homogeneous part of $E^*_\mu$ is 
    %the nonsymmetric Macdonald polynomial 
    $E_\mu$.
\end{thm}



%The definition of these polynomials will not be needed here, but we will recall some of their useful properties.

When indexed by a partition $\lambda$, nonsymmetric interpolation Macdonald polynomials coincide with interpolation ASEP polynomials: $F^*_\lambda=E^*_\lambda$. We also have the following 
result from \cite[Lemma 2.5 and Proposition 2.14]{BenDaliWilliams2025}.
\begin{prop}[\cite{BenDaliWilliams2025}]\label{prop:Vstar}
    Fix  a partition $\lambda$. Then $\{E_{\mu}^* \ \vert \ \mu\in S_n(\lambda)\}$ and $\{F_{\mu}^* \ \vert \ \mu\in S_n(\lambda)\}$ are both bases for $\VV$. In particular, if $\mu\in S_n(\lambda)$, then $F^*_\mu$ is a linear combination of $E_{\nu}^*$ for $\nu\in S_n(\lambda)$.
\end{prop}
\begin{remark}\label{rmk:basis_1}
In the following, we will use the fact that even after evaluating at $q=1$, $\{E^*_\mu(\bfx;1,t) \ \vert \ {\mu\in \YY_n}\}$ and $\{F^*_\mu(\bfx;1,t) \ \vert \ {\mu\in \YY_n}\}$  are both bases of the polynomial space in $n$ variables and coefficients in $\QQ(q)$. Indeed, $E^*_\mu$ is triangular in the monomial basis (with respect to the dominance order on compositions) (see \cite[Theorem 3.11]{Knop1997b}) and the coefficient of $x_\mu$ in $E^*_\mu$ is 1, independent of $q$. A similar argument works for the polynomials $F^*_\mu$ using \cref{eq:normalization_F}.
\end{remark}

    

We will make use of the \emph{dehomogenization map}; see \cite[Definition 3.1 and Theorem 3.8]{NakviSahiSergel2023} for a similar definition in the context of Jack polynomials.
\begin{definition}\label{def:dehom}
We define a linear map $\boldsymbol{\psi}:\mcP_n \to \mcP_n$ 
%of polynomials in $n$ variables 
by 
\begin{equation}
 \boldsymbol{\psi}: E_\mu\mapsto E_\mu^*, \quad\text{for any partition $\mu$.}   
\end{equation}
We call the operator $\boldsymbol{\psi}$ the \emph{dehomogenization map}.
\end{definition}

Since both $(E_\mu)$ and $(E_\mu^*)$ are bases of $\mcP_n$, this map is well defined and is an isomorphism.
Alternatively, $\boldsymbol{\psi}$ can be defined as the linear operator such that if $f$ is homogeneous of degree $d$ then $\boldsymbol{\psi}(f)$ is the unique polynomial of top homogeneous part $f$ and which vanishes on all compositions of size smaller than $d$. 

Using %the vanishing conditions of the polynomials $F^*_\mu$ from 
\cref{prop:Vstar}, and the fact that the top homogeneous part of $F^*_\mu$ is $F_\mu$, we get that
$\boldsymbol{\psi}(F_\mu)=F_\mu^*$. Thus we can conclude the following.

\begin{lem}\label{lem:dehomog}
If we have an expansion 
$E_\mu=\sum_{\nu}c_{\mu,\nu}F_\nu$
(for some coefficients $c_{\mu,\nu}\in \QQ(q,t)$), we get that 
$E^*_\mu=\sum_{\nu}c_{\mu,\nu}F^*_\nu$ by applying the map $\boldsymbol{\psi}$.
\end{lem}

\subsection{The Hecke operators}
We now recall the definition of the Hecke operators and their action on the ASEP polynomials (see e.g. \cite[Section 2.3]{BenDaliWilliams2025}).

For $1\leq i \leq n-1$, we let $s_i=(i,i+1)$ denote the transposition exchanging $i$ and $i+1$. The transposition $s_i$ acts on a polynomial in $\mcP_n$ by permuting the variables $x_i$ and $x_{i+1}$:
$$s_i\cdotp f(x_1,\dots,x_i,x_{i+1},\dots,x_n)
=f(x_1,\dots,x_{i+1},x_i,\dots,x_n).$$

The \emph{Hecke operator} $T_i$ is the linear operator on $\mcP_n$ defined by 
\begin{equation}\label{eq:Hecke}
T_i:=t-\frac{tx_i-x_{i+1}}{x_i-x_{i+1}}(1-s_i).
\end{equation}

These operators satisfy the relations of the Hecke algebra of type $A_{n-1}$
\begin{equation}\label{eq:Hecke_algebra}
\begin{array}{ll}
(T_i-t)(T_i+1)=0 & \text{for $1\leq i \leq n-1$}     \\
 T_iT_{i+1}T_i=T_{i+1}T_iT_{i+1}    & \text{for $1\leq i \leq n-2$}\\
 T_iT_j=T_jT_i    & \text{for $ |i-j| > 1$}.\\
\end{array}
\end{equation}

%\LW{added this} The following lemma is immediate.
%\begin{lem}\label{lem:Heckesymm}
%If $f$ is a polynomial which is symmetric in $x_i$ and $x_{i+1}$, then $T_i(f) = tf$.
%\end{lem}

If $\sigma\in S_n$ and $\sigma=s_{i_1}\dots s_{i_\ell}$ is a reduced decomposition of $\sigma$, we define 
\begin{equation}\label{def:Hecke}
T_\sigma:=T_{i_1}\dots T_{i_\ell}.
\end{equation}
It follows from \eqref{eq:Hecke_algebra} that this definition is independent of the choice of reduced expression. 


\begin{prop}\label{prop:T_i-F_star} \cite[Proposition 2.10]{BenDaliWilliams2025}
    Let $\mu = (\mu_1,\dots,\mu_n)$.  For $1\leq i \leq n-1$, the interpolation ASEP polynomials $F^*_\mu$ satisfy the following: % properties:
    \begin{enumerate}
        \item\label{item 1} $T_i F^*_{\mu}=F^*_{s_i\mu}$ if $\mu_i>\mu_{i+1}$,
        \item\label{item 2} $T_iF^*_{\mu}=tF^*_{\mu}$ if $\mu_i=\mu_{i+1}$,
        \item\label{item 3} $T_i F^*_{\mu}=(t-1)F^*_{\mu}+tF^*_{s_i\mu}$ if $\mu_i<\mu_{i+1}$.
    \end{enumerate}
\end{prop}

Moreover, we have the following \emph{shape permuting
operator} formula for interpolation nonsymmetric Macdonald polynomials.

\begin{prop}[{\cite[Proposition 3.6]{BenDaliWilliams2025}}]
\label{prop:shape}
Let $\nu\in \NN^n$ be a composition and fix $1\leq j \leq n-1$. Suppose that $\nu_j>\nu_{j+1}$, and
write
\begin{equation}\label{eq:def_r}
r_j(\nu):=
\#\{k<j \mid \nu_{j+1}<\nu_{k}\leq \nu_j\}
+
\#\{k>j \mid \nu_{j+1}\leq \nu_k < \nu_j\}.
\end{equation}

Then
    \begin{equation}
\label{shape-permute}
E^*_{s_j \nu}(\bfx;q,t) = 
\left(T_j + \frac{1-t}{1-q^{\nu_{j}-\nu_{j+1}}t^{r_j(\nu)}}\right)
E^*_{\nu}(\bfx;q,t).
\end{equation}
\end{prop}

%Knop and Sahi have proved the following recurrence for nonsymmetric interpolation Macdonald polynomials (see \cite[Corollary 2.5]{Knop1997b}).

%\HBD{Do we need the Knop–Sahi recurrence?} \LW{Now I guess we keep it because of the modified proof in the appendix}\HBD{Yes!}
\begin{thm}[{\cite{Knop1997b,Sahi1996}}]\label{thm:KS_recurrence}
    For any composition $\mu$, we have
 \begin{equation}\label{eq:KS_recurrence}
  E^*_{(\mu_n+1,\mu_1,\dots,\mu_{n-1})}(x_1,\dots,x_n;q,t)=\left(x_1-\frac{1}{t^{n-1}}\right)q^{\mu_n}
  E_\mu^*\left(x_2,\dots,x_n,\frac{x_1}{q};q,t\right).  
    \end{equation}
\end{thm}

\begin{comment}
\begin{remark}
One can check that \cref{eq:KS_recurrence} and the special case of \cref{shape-permute} 
% \LW{probably mean \cref{shape-permute}?} 
in which $\nu_{j+1}=0$  together with the base case $E^*_{(0,\dots,0)}=1$ uniquely determine the polynomials $E_\mu^*$. See \cite[Lemma 2.1.2]{HaglundHaimanLoehr2008} for a similar argument in the homogeneous case.
\end{remark}
\end{comment}

\section{The proof of \texorpdfstring{\cref{thm:weak-reordering}}{}}\label{app:permuted_basement}
%\subsection{Definition}
In this appendix we use 
% we define an interpolation analogue $E_\mu^{*\sigma}$ of the permuted-basement polynomials $E_\mu^\sigma$ introduced in \cite{ferreira-2011,Alexandersson2019}.
% We then use them, together with 
analogues of results from \cite{AlexanderssonSawhney2019,AyyerMartinWilliams2025} to 
prove \cref{thm:weak-reordering}.

%The convention we use here to 

% \begin{remark}
%  We will define $E_\mu^{*\sigma}$  so that its top homogeneous part gives $E_\mu^\sigma$ as defined in \cite{Alexandersson2019}.
%     Regarding conventions, 
%     our notation $E_{\mu}^{\sigma}$ agrees with the notation $E_{\mu}^{\sigma}$ 
% from \cite{Alexandersson2019} and \cite{AlexanderssonSawhney2019}.  Our notation $E_{\mu}$ agrees with the 
% notation $E_{\mu}$ from \cite{HaglundHaimanLoehr2008} and the notation
% $E_{{\mu}^{\rev}}$ from \cite{Alexandersson2019} and \cite{AlexanderssonSawhney2019}, 
% where $\mu^{\rev}=(\mu_n,\dots,\mu_1)=w_0\cdotp \mu,$ and
% $w_0$ is the longest element in $S_n$.  In one-line notation, we have 
% $w_0:=(n,n-1,\dots,1).$
% \end{remark}

    
    
%\LW{Is our convention definitely the same as in 
%\cite{Alexandersson2019} and \cite{AlexanderssonSawhney2019}? When I compare \eqref{eq:reverse} to page 6 of \cite{AlexanderssonSawhney2019} I'm not so sure.  Just below Def 5 on page 6, they suggest that $E_{\lambda}^{w_0} = E_{\lambda}$, though it's not totally clear what conventions they are using here since they cite HHL.  At any rate, we should state very clearly at the start of this section whether we agree with the AS conventions or if there is a reversal.}


% \begin{definition}\label{def:permuted_basement}
% Let $\sigma$ be a permutation in $S_n$. Define the 
% \emph{interpolation permuted-basement polynomial} $\Estar{\sigma}{\mu}$ by induction on $\sigma$ as follows.
% When $\sigma=w_0$, we define
% \begin{equation}\label{eq:reverse}
% \Estar{w_0}{\mu}:=E^*_{\mu^{\rev}}, \text{ where }
% %\end{equation}
% %where 
% %$$
% \mu^{\rev}=(\mu_n,\dots,\mu_1)=w_0\cdotp \mu.
% \end{equation}

% And for $1\leq i\leq n-1$ such that $\ell(s_i\sigma)<\ell(\sigma)$, we define
% \LW{I turned around the equation below, since the LHS used to have $T_i$ in it.}
% \begin{equation}\label{eq:basement_permuting}
% \Estar{s_i\sigma}{\mu}:=
%   T_i \Estar{\sigma}{\mu}
%     \times\begin{cases}
%      t^{-1}  &\text{ if } \mu_{\sigma_i^{-1}}\leq\mu_{\sigma_{i+1}^{-1}}\\
%      1&\text{ otherwise}.
% \end{cases}
% \end{equation}

% % \LW{How about phrasing it like this instead?
% % \begin{equation}\label{eq:basement_permuting}
% % \Estar{s_i\sigma}{\mu}:=
% %   T_i \Estar{\sigma}{\mu}
% %     \times\begin{cases}
% %      t^{-1}  &\text{ if } \mu_{\sigma_i^{-1}}\leq\mu_{\sigma_{i+1}^{-1}}\\
% %      1&\text{ otherwise}.
% % \end{cases}
% % \end{equation}
% % }
% % \HBD{In \cite{AlexanderssonSawhney2019} it is $\sigma s_i$ instead of $s_i\circ \sigma$. Do they use another convention to compose permutations?}

% Equivalently, if $\sigma=\tau \ w_0 $
% % , and $\tau=s_{i_1}s_{i_2}\dots s_{i_\ell}$ is a reduced expression
% , then
% \begin{equation}\label{eq:Estar}
% \Estar{\sigma}{\mu}:=T_\tau \ \Estar{w_0}{\mu}\times t^{-r},
% \end{equation}
% where $T_\sigma$ is the operator defined by \cref{def:Hecke} and
% \begin{equation}\label{eq:r}
% % T_\tau:=T_{i_1}T_{i_2}\dots T_{i_\ell} \text{ and }
%   r:=\#\left\{i<j:\sigma_i< \sigma_j\text{ and }
%   \mu_i\geq \mu_j\right\}.
% \end{equation}
    
% \end{definition}


% \begin{example}
%     Let $n=4$, $\mu=(3,4,0,4)$, and $\sigma=(3,1,2,4)=(4,2,1,3)w_0.$ We have the reduced expression
%     $(4,2,1,3)=s_3 s_1 s_2 s_1.$
%     Then 
%     $$\Estar{\sigma}{\mu}=
%     \Estar{s_3 s_1 s_2 s_1 w_0}{(3,4,0,4)}=
%     t^{-2}T_3 T_1 T_2 T_1 \Estar{w_0}{(3,4,0,4)}= 
%     t^{-2}T_3 T_1 T_2 T_1 \Estar{}{(4,0,4,3)}.$$
%     Here 
%     the pairs $(i,j)$ contributing to \cref{eq:r} are $(2,3)$ and $(2,4)$ so $r=2$.
% \end{example}

% \LW{I added lemma below}
% \begin{lem}
%      The top homogeneous part of $E_\mu^{*\sigma}$  equals $E_\mu^\sigma$ as defined in \cite{Alexandersson2019}.
% \end{lem}
% \begin{proof}
%     This follows from \cite[Proposition 7]{AlexanderssonSawhney2019} and 
%     \cite[Section 3 and footnote 1]{Alexandersson2019}.
% \end{proof}


% \begin{remark}
%     Permuted-basement non-symmetric Macdonald polynomials 
%     were defined in \cite{ferreira-2011} combinatorially, see also 
%     %have been defined combinatorially in 
%     \cite{Alexandersson2019}. 
%     However here we define their interpolation analogue using the action of the Hecke operators.
%     The operators $\widetilde{\theta}_i$ and $\widetilde{\pi}_i$ in \cite{AlexanderssonSawhney2019}, correspond respectively to $T_i$ and  $T_i+1-t$ in our notation.
% \end{remark}


\begin{comment}
\HBD{It seems that we don't need this technical corollary after all. In \cite{AlexanderssonSawhney2019}, it's used to prove the homogeneous version of \cref{lem:E_one_row}. But, we can derive the inhomogeneous one by applying the map $\boldsymbol{\psi}$.}


The following is an analogue of \cite[Proposition 8]{AlexanderssonSawhney2019}. The proof is an algebraic version of the proof of \cite[Proposition 17]{Alexandersson2019}.
\LW{Take a look at this soon.}
\begin{cor}
    Let $\nu$ be a composition and $\sigma$ a permutation. Suppose $\nu_j<\nu_{j+1}$, $\sigma(j)=i+1$ and $\sigma(j+1)=i$ for some $i$ and $j$.
Write
\[
r'_j(\nu)=
\#\{k>j+1 \mid \nu_{j}<\nu_{k}\leq \nu_{j+1}\}
+
\#\{k\leq j \mid \nu_{j}\leq \nu_k < \nu_{j+1}\}.
\]
Then, for any permutation $\sigma$, we have
    \begin{equation}
\label{eq:shape-permute}
\Estar{\sigma}{s_j \nu}(\bfx;q,t) = 
\left(T_i + \frac{1-t}{1-q^{\nu_{j+1}-\nu_j}t^{r'_j(\nu)}}\right)
\Estar{\sigma}{\nu}(\bfx;q,t).
\end{equation}
\end{cor}

Note that the statistics $r'$ associated to a composition $\nu$, correspond to the statistics $r$ associated to $\nu^{\rev}$ as defined in \cref{eq:def_r}. 
More precisely, $r'_j(\nu)=r_{n-j}(\nu^{\rev})$ for any $1\leq j\leq n-1$.
This reversal of convention compared to \cref{eq:shape-permute} is related to the reverse appearing in the definition of $E_\mu^{*\sigma}$ \cref{def:permuted_basement}.
\begin{proof}
    When $\sigma=w_0$, the result corresponds to \cref{prop:shape}.
    We now fix $\nu$ and $j$, we want to prove that \cref{eq:shape-permute} holds for all pairs $(\sigma,i)$ such that $\sigma(j)=i+1$ and $\sigma(j+1)=i$. We will show the following:
    \begin{itemize}
        \item Fix $k\notin\{i-1,i,i+1\}$. Then \cref{eq:shape-permute} holds for $(\sigma,i)$ if and only if it holds for $(s_k\sigma,i)$.
        \item Assume $i\geq 2$ and that 
        \cref{eq:shape-permute} holds for $(\sigma,i)$, then it also holds for $(s_i s_{i-1}\sigma,i-1)$.
        \item Assume $i\leq n-2$ and that 
        \cref{eq:shape-permute} holds for $(\sigma,i)$, then it also holds for $(s_i s_{i+1}\sigma,i+1)$.
        % \item If
        % \cref{eq:shape-permute} holds for $(\sigma,i,j)$ then it holds for $(s_i s_{i-1}\sigma,i-1,j)$.
    \end{itemize}
    \HBD{The second and the third items are more general operations than the ones given in \cite{Alexandersson2019}.}
    
    One can check that any pair $(\sigma,i)$ satisfying the conditions of the corollary can be obtained from $(w_0,n-j+1)$ using the three operations above.
    
    % \HBD{For any pair $(\sigma,i)$ as above, $\sigma$ can be written
    % $$\sigma=s_{k_1}\dots s_{k_m}w_0,$$
    % and such that for any
    % $1\leq u\leq m$, we have on of the following
    % \begin{itemize}
    %     \item $k_u\notin\{i-1,i,i+1\}$,
    %     \item $k_u=i$ and $k_{u+1}\in\{i-1,i+1\}$,
    %     \item $k_u\in\{i-1,i+1\}$ and $k_{u-1}=i$,
    % \end{itemize}
    % }

    
    We prove the first item. Let $k\notin\{i-1,i,i+1\}$. We distinguish the two cases $\ell(s_k\sigma)<\ell(\sigma)$ and $\ell(s_k\sigma)>\ell(\sigma)$.
    
    \textbf{The case $\ell(s_k\sigma)<\ell(\sigma)$:}
    By definition we have,
    $$C\times\Estar{s_k \sigma}{s_j\nu}\ =T_k\Estar{\sigma}{s_j\nu}.$$
where $$C:=
\begin{cases}
     t  &\text{ if } (s_j\nu)_{\sigma^{-1}(k)}\leq (s_j\nu)_{\sigma^{-1}(k+1)}\\
     1&\text{ otherwise}.
\end{cases}
$$
Similarly
\begin{align*}
  C\times\left(T_i + \frac{1-t}{1-q^{\nu_{j+1}-\nu_j}t^{r'_j(\nu)}}\right)
\Estar{s_k\sigma}{\nu}(\bfx;q,t)
&=\left(T_i + \frac{1-t}{1-q^{\nu_{j+1}-\nu_j}t^{r'_j(\nu)}}\right)
T_k\Estar{\sigma}{\nu}(\bfx;q,t)\\
&=T_k\ \left(T_i + \frac{1-t}{1-q^{\nu_{j+1}-\nu_j}t^{r'_j(\nu)}}\right)
\Estar{\sigma}{\nu}(\bfx;q,t),
\end{align*}
where we use the fact that $T_k$ and $T_i$ commute (since $k\neq i-1,i+1$).
Now we use the fact the result holds for the pair $(\sigma,i)$ to get

$$C\times \left(T_i + \frac{1-t}{1-q^{\nu_{j+1}-\nu_j}t^{r'_j(\nu)}}\right)
\Estar{s_k\sigma}{\nu}(\bfx;q,t)=T_k\cdotp E_{s_j\nu}^{*\sigma}=C'\times E_{s_j\nu}^{*s_k\sigma},$$
where 
$$C'
=\begin{cases}
     t  &\text{ if } \nu_{\sigma^{-1}(k)}\leq \nu_{\sigma^{-1}(k+1)}\\
     1&\text{ otherwise}.
\end{cases}$$
It remains to check that $C'=C$. Indeed, $s_j\nu$ is obtained from $\nu$ by permuting entries $j$ and $j+1$ but $\sigma^{-1}(k),\sigma^{-1}(k+1)\notin\{j,j+1\}$. This proves that $C'=C$ and as a consequence that 
$$\left(T_i + \frac{1-t}{1-q^{\nu_{j+1}-\nu_j}t^{r'_j(\nu)}}\right)
\Estar{s_k\sigma}{\nu}(\bfx;q,t)=E_{s_j\nu}^{*s_k\sigma},$$
as desired.

\textbf{The case $\ell(s_k\sigma)>\ell(\sigma)$}. We can write
    $$C^{-1}\times\Estar{s_k \sigma}{s_j\nu}\ =T_k^{-1}\cdotp\Estar{\sigma}{s_j\nu},$$
 
where $$C:=
\begin{cases}
     t  &\text{ if } (s_j\nu)_{\sigma^{-1}(k+1)}\leq (s_j\nu)_{\sigma^{-1}(k+1)}\\
     1&\text{ otherwise},
\end{cases}
$$
we then proceed as above.
This proves the first item of the corollary.

We obtain in a similar way the two other items, by applying $T_iT_{i-1}$ and $T_iT_{i+1}$ respectively, and using the fact that Hecke operators satisfy the braid relations.
\end{proof}

% \HBD{In the following two subsections, most of the proofs work in the same way as in the homogeneous case, with the (slight) exception of \cref{lem:E_one_row} and  \cref{prop:add_one}.}

\end{comment}

% \subsection{Interpolation analogues of some results from \cite{AlexanderssonSawhney2019}}

%\HBD{Should we remove this paragraph?}
% In this subsection, we state the interpolation analogues of some results from \cite{AlexanderssonSawhney2019}. 
%Most of these results are obtained in the same way as in the homogeneous case. This is mainly because they are based on the shape permuting formula 
%(\cref{prop:shape})
%and Knop–Sahi recurrence
%(\cref{thm:KS_recurrence}), 
%both of which are satisfied by the interpolation polynomials.


% \begin{lem}\label{lem:E_one_row}
% For any $0\leq m\leq n$, we have
% $$\Estar{\sigma}{1^m,0^{n-m}}(x_1,\dots,x_n;1,t)=e_m^*(x_1,\dots,x_n;t).$$
% \end{lem}

% \begin{proof}
% We will start by proving the lemma when $\sigma=w_0;$ that is, we will first prove that 
% \begin{equation}\label{eq:first}
%    E^*_{0^{n-m},1^m}(x_1,\dots,x_n;1,t)=e_m^*(x_1,\dots,x_n;t).
% \end{equation}
% Let us expand the non-symmetric homogeneous polynomial $E_{0^{n-m},1^m}$ in the ASEP basis:
% \begin{equation}\label{eq:E_f}
%   E_{0^{n-m},1^m}(x_1,\dots,x_n;q,t)=\sum_{\mu\in S_n(1^m,0^{n-m})}d_{\mu}(q,t)F_{\mu}(x_1,\dots,x_n;q,t),  
% \end{equation}
% for some coefficients $d_\mu(q,t)$. Recall that for $\mu\in S_n(1^m,0^{n-m})$, $F_\mu$ is independent of $q$ (see \cref{eq:factorization_one_row}) \HBD{We are not really using this, right?}.
% Moreover, we know from \cref{eq:reverse} and  \cite[Lemma 13]{AlexanderssonSawhney2019} %\LW{this lemma says that $E^{\sigma}_{1^m 0^n} = e_m$; so is there a reversal in the indexing of compositions between them and us? or are we using a statement like when $\sigma=w_0$, we get a reversal of the composition?}
% %\HBD{Yes, we are using their lemma for $\sigma=w_0$, for which we have from \cref{eq:reverse} that $E_{0^{n-m},1^m}=E^{w_0}_{1^m,0^{n-m}}=e_m$}
% that at $q=1$, 
% \begin{equation}
%   E_{0^{n-m},1^m}(x_1,\dots,x_n;1,t)=E^{w_0}_{1^m,0^{n-m}}=e_m(x_1,\dots,x_n;t).  
% \end{equation}
% Combining this with \cref{eq:P_column_hom}, and using the fact that $(F_\mu(q=1))_{\mu\in \YY_n}$ is a basis (see \cref{rmk:basis_1} we get that $d_\mu(1,t)=1$.

% By applying the dehomogenization map $\boldsymbol{\psi}$ to \cref{eq:E_f} and using \cref{lem:dehomog}, we get
% \begin{equation}\label{eq:expansion}
% E^*_{0^{n-m},1^m}(x_1,\dots,x_n;q,t)=\sum_{\mu\in S_n(1^m,0^{n-m})}d_\mu(q,t)F^*_{\mu}(x_1,\dots,x_n;q,t).
% \end{equation}
% We now make the substitution $q=1$.
% Using \cref{eq:expansion} and \cref{eq:P_column}, we get
% $$E^*_{0^{n-m},1^m}(x_1,\dots,x_n;1,t)=\sum_{\mu\in S_n(1^m,0^{n-m})}F^*_{\mu}(x_1,\dots,x_n;q,t)=e_m^*(x_1,\dots,x_n;1,t).$$
% This gives \cref{eq:first}, and hence proves the lemma when $\sigma=w_0$.

%     % We know from \cref{prop:shape} that for any $0\leq j\leq m-1$ and $0\leq k\leq n-m-1$ that
%     % $$E^*_{1^j,0^{k+1},1,0^{n-m-k-1},1^{m-j}}=\left(T_{j+k}+\frac{1-t}{1-t^{n-m-k+j}}\right)E^*_{1^j,0^{k},1,0^{n-m-k},1^{m-j}}.$$
%     % But we know that $E^*_{1^m,0^{n-m}}=F^*_{1^m,0^{n-m}}.$ We can prove by induction 
% We now obtain the formula for $\Estar{\sigma}{1^m,0^{n-m}}$ by 
% using \cref{eq:Estar} and \cref{eq:first}. %to   
% %multiple application of the Hecke operators on 
% $\Estar{w_0}{1^{n-m},0^m}$. % using \cref{eq:basement_permuting}.
% We also use \cref{lem:Heckesymm} and the fact that $e_m^*$ is symmetric.
% %, and that Hecke operators act on a symmetric function as multiplication by $t$.
% \end{proof}

% Fix a partition $\lambda$. The \emph{dominance order} 
% on $S_n(\lambda)$ is the partial order defined by   
% $\mu\leq \nu$ if and only if
% $$\mu_1+\dots+\mu_i\leq \nu_1+\dots+\nu_i,\quad \text{for any $1\leq i\leq n$}.$$


% \begin{definition}
%     The \emph{weak standardization} of a composition $\mu$ is the smallest composition $\WS{\mu}$ in the lexicographic order such that $\WS{\mu}_{i}\leq \WS{\mu}_j$ if and only if $\mu_{i}\leq \mu_j$ for all pairs $i,j$. In other terms, $\WS{\mu}$ is obtained from $\mu$ by replacing the smallest entries by 0, the second smallest entries by 1, and so on.
% \end{definition} 

% \begin{example}
%     If $\mu=(2,1,2,4)$, then $\WS{\mu}=(1,0,1,2).$
% \end{example}

% The following is an analogue of 
% \cite[Lemma 14]{AlexanderssonSawhney2019}.
% \begin{lem}\label{lem_Bruhat_induction}
%     Let $\mu,\nu\in S_n(\lambda)$, such that $\mu\leq \nu$. We assume that 
%     $\frac{\Estar{w_0}{\mu}(\bfx;1,t)}{\Estar{w_0}{\WS{\mu}
%     }(\bfx;1,t)}$ is symmetric. Then 
%     $$\frac{\Estar{w_0}{\mu}(\bfx;1,t)}{\Estar{w_0}{\WS{\mu}
%     }(\bfx;1,t)}=\frac{\Estar{w_0}{\nu}(\bfx;1,t)}{\Estar{w_0}{\WS{\nu}
%     }(\bfx;1,t)}.$$
% \end{lem}

% \begin{proof}
%     We can go from $\Estar{w_0}{\mu}(\bfx;1,t)$ to $\Estar{w_0}{\nu}(\bfx;1,t)$ by applying the shape-permuting operators
%     from \cref{prop:shape}. At $q=1$, \LW{I added more detail here but dropped reference \cite{AlexanderssonSawhney2019}} these operators have the same behavior on both $\Estar{w_0}{\mu}$ and 
%     $\Estar{w_0}{\WS{\mu}}$ because $r_j(\mu) = r_j(\WS{\mu})$.
%     %are invariant under taking the weak standardization; see 
%     %See  \cite{AlexanderssonSawhney2019} for more details.
% \end{proof}

% The following is an analogue of \cite[Equation 10]{AlexanderssonSawhney2019}.
% \begin{prop}\label{prop:add_one}
%     For any composition $\mu=(\mu_1,\dots,\mu_n)$ and any permutation $\sigma\in S_n$, we have
%     \begin{equation}
% \Estar{\sigma}{\mu+\mathbf{1}}(x_1,\dots,x_n; q, t)=\left(x_1-\frac{1}{t^{n-1}}\right)\dots \left(x_n-\frac{1}{t^{n-1}}\right)q^{|\mu|}\Estar{\sigma}{\mu}\left(\frac{x_1}{q},\dots,\frac{x_n}{q};q,t\right),
%     \end{equation}
% where $\mu+\mathbf{1}=(\mu_1+1,\dots,\mu_n+1)$.
% \end{prop}
% \begin{proof}
% We start by proving the case $\sigma=w_0$; by \cref{eq:reverse}, 
% it suffices to show that for any $\mu$ we have
% \begin{equation}\label{eq:add_one}
% \Estar{}{\mu+\mathbf{1}}(x_1,\dots,x_n; q, t)=\left(x_1-\frac{1}{t^{n-1}}\right)\dots \left(x_n-\frac{1}{t^{n-1}}\right)q^{|\mu|}\Estar{}{\mu}\left(\frac{x_1}{q},\dots,\frac{x_n}{q};q,t\right).
%     \end{equation}
% Let $g$ denote the right-hand side of \cref{eq:add_one}. We check that $g$ satisfies the characterizing properties of $E^*_{\mu+\mathbf{1}}$ from \cref{thm:Knop_Sahi}. %_1+1\dots \mu_n+1}$. 
% First, we have that 
%     $$\left[\prod_ix_i^{\mu_i+1}\right]g=[x^\mu]E^*_\mu=1.$$
% Fix now a partition $\nu$ such that $|\nu|\leq|\mu+\mathbf{1}|=|\mu|+n$ and $\nu\neq \mu+\mathbf{1}$. We claim that $g(\widebar{\nu})=0$.
% % \LW{I think we want $g(\widebar{\nu})=0$} 
% We distinguish two cases:
% \begin{itemize}
%     \item If there exists $i$ such that $\nu_i=0$, then we take the minimal such $i$ and we have $\widebar{\nu}_i=t^{-n+1}.$
%     Thus, when we evaluate $g$ at $\widebar{\nu}$, 
%     $(x_i-\frac{1}{t^{n-1}})$ vanishes, so the claim follows.
%     \item If $\nu_i>0$ for $1\leq i\leq n$, then $\widebar{\nu}=q\times\widebar{\kappa}$, where $\kappa=(\nu_1-1,\dots,\nu_n-1)$. Since $|\kappa|\leq|\mu|$ and $\kappa\neq \mu$, then $E^*_\mu$ vanishes at $\widebar{\kappa}$, and we deduce that
%     $E_\mu^*(x_1/q,\dots,x_n/q; q, t)$ vanishes at $\widebar{\nu}$.
% \end{itemize}
% This gives the proposition when $\sigma=w_0$. We then get any other permutation $\sigma$ by applying Hecke operators; see \cref{eq:basement_permuting}. We use here both \cref{lem:Heckesymm} and the fact that $\left(x_1-\frac{1}{t^{n-1}}\right)\dots \left(x_n-\frac{1}{t^{n-1}}\right)$ is symmetric.
% \end{proof}

% \HBD{I think that the proofs of 
% \cite[Lemma 15, Proposition 16 and Theorems 17 and 18]{AlexanderssonSawhney2019} remain exactly the same.}

% The following is an analogue of \cite[Theorem 17]{AlexanderssonSawhney2019}. 
% It is obtained by induction on $\sigma$ and $\lambda$ using \cref{lem_Bruhat_induction}, \cref{prop:add_one} and the basement-permuting formula \cref{eq:basement_permuting}.  

%\LW{Wouldn't we also need to use here \cref{thm:KS_recurrence}?  Otherwise I'm not sure how we could get shapes besides say 
%$(0,0,0)$, %$(1,1,1)$, $(2,2,2)$, etc.
%And maybe we should say more on the strategy of proof.  Perhaps that the 
%base case is $%\mu=(0,0,\dots,0)$, and we first use induction on the size of $\mu$ to prove it for $\Estar{w_0}{\mu}$?}
%\HBD{I removed the paragraph before the thm and added a short proof.}
%\LW{The theorem below is actually more specific than
%\cite[Theorem 17]{AlexanderssonSawhney2019}.  Maybe we should phrase this one similarly, and say that the RHS is a ratio of symmetric polynomials (with coefficients in $\QQ(t)$)?  
%Is the idea to prove that first, then use Theorem B.10 to 
%pin down that ratio of symmetric polynomials?}
%\LW{I don't find the proof sketches of Thm B.9 and B.10 terribly helpful; what is the base case for the inductions?}

% \HBD{Added these definitions and the convention about $e_0^*$ in \cref{def:int_elem}.}
% \LW{Great: I moved the defn of $e_{\lambda}^*$ into Def 4.1}
If $\lambda$ is a partition, we will denote by $\lambda'$ \textit{the conjugate} of $\lambda$.
Recall that $e_{\lambda}^*$ is the interpolation elementary symmetric function defined in \cref{def:int_elem}. Moreover, if $\mu=(\mu_1,\dots,\mu_n)$ is a composition, we denote $\mu^{\rev}:=(\mu_n,\dots,\mu_1)$
% \begin{thm}\label{thm:standardization}
% Fix a composition $\mu$ and permutation $\sigma$. Let $\lambda$ be the partition obtained by reordering $\mu$.
%     We have
%         $$\frac{\Estar{\sigma}{\mu}(x_1,\dots,x_n;1,t)}{\Estar{\sigma}{\widetilde{\mu}}(x_1,\dots,x_n;1,t)}=\frac{e^*_{\lambda'}(x_1,\dots,x_n;t)}{e^*_{\widetilde{\lambda}'}(x_1,\dots,x_n;t)}.$$
% \end{thm}
% We give the main ideas of the proof of this theorem. We refer to the proof of \cite[Theorem 17]{AlexanderssonSawhney2019} for more details.
% \begin{proof}
%     % It is enough to prove the theorem for $\sigma=w_0$, and the general case will follow from the basement-permuting formula \cref{eq:basement_permuting}.
%     We proceed by induction on $\mu$, first with respect to its size $|\mu|$, and, among compositions of the same size, with respect to the lexicographic order.
%     If $\mu=(0,\dots,0)$, then both sides are equal to 1. If $\mu_i\geq 1$ for all $i$, we apply \cref{prop:add_one} to decrease the size of the composition. Using \cref{lem_Bruhat_induction}, we can reduce to the case of a partition, in which case we can apply \cref{thm:Knop_Sahi}. 
%     \HBD{This last step is not obvious}.
% \end{proof}


\begin{thm}[{\cite[Theorem 18]{AlexanderssonSawhney2019}}]
\label{thm:E_partition_hom}
    For any partition $\lambda$, we have
    $$E^{}_{\lambda^{\rev}}(x_1,\dots,x_n;1,t)=e_{\lambda'}(x_1,\dots,x_n).$$
\end{thm}

We now give  a dehomogenization of this theorem.
% It can be obtained by induction from \cref{lem:E_one_row}, the shape-permuting formula \cref{shape-permute} and the Knop–Sahi recurrence \cref{thm:KS_recurrence}.
% \LW{I am confused on how this proof would go, because the shape-permuting formula and Knop-Sahi recurrence will take us out of the world of partitions.  I wonder if we could prove the result below using the dehomogenization operator?} \LW{We could mention that we prove it first in the case $\sigma=w_0$ and then use basement-permuting formula.}
% \HBD{I added an independent proof. 
% It is almost the same as the proof of \cref{lem:E_one_row}.}

\begin{thm}\label{thm:E_partition}
    For any partition $\lambda$, we have
    $$E^*_{\lambda^{\rev}}(x_1,\dots,x_n;1,t)=e^*_{\lambda'}(x_1,\dots,x_n;t).$$
\end{thm}
\begin{proof}
% It is enough to prove the theorem for $\sigma=w_0$, and the general case will follow from the basement-permuting formula \cref{eq:basement_permuting}.
From \cref{prop:Vstar}, we can write an expansion
$$E_{\lambda^{\rev}}(\bfx;q,t)=\sum_{\mu\in S_n(\lambda)}d_{\mu}(q,t)F_\mu(\bfx;q,t),$$
for some coefficients $d_\mu\in \QQ(q,t)$.
Applying the dehomogenization map \cref{lem:dehomog}, we get
\begin{equation}\label{eq:E_F}
  E^*_{\lambda^{\rev}}(\bfx;q,t)=\sum_{\mu\in S_n(\lambda)}d_{\mu}(q,t)F^*_\mu(\bfx;q,t).  
\end{equation}
    But we know from \cref{prop:symmetrization} and \cref{thm:factorization} that
    % \LW{I think the next two equations should have a $*$ on the LHS}
    $$e^*_{\lambda'}(\bfx;t)=\sum_{\mu\in S_n(\lambda)}F^*_\mu(\bfx;1,t).$$
    Combining this with  \cref{thm:E_partition},%(applied for $\sigma=w_0$), 
        we get
    $$E^*_{\lambda^{\rev}}(\bfx;1,t)=\sum_{\mu\in S_n(\lambda)}F^*_\mu(\bfx;1,t).$$
    Comparing this with \cref{eq:E_F}, and the fact that, interpolation ASEP $F_\mu^*$ polynomials are a basis (see \cref{rmk:basis_1}), we deduce that $d_\mu(1,t)=1$ for any $\mu\in S_n(\lambda)$ as desired.
\end{proof}

% \subsection{Interpolation analogues of some results from \cite{AyyerMartinWilliams2025}}


% The following is an analogue of \cite[Theorem 4.1]{AyyerMartinWilliams2025}.
% The proof is the same.

% \begin{thm}\label{thm:adjacent}
%     Consider a composition $\eta$ such that its parts of sizes $a$ and
% $b$ (with $a<b$) occur in increasing order from left to right in $\eta$, and such that $\eta$ has no parts of size $h$ for any $a < h < b$. (In this case, we say that the classes $a$ and $b$ are
% adjacent.) Let $\widehat{\eta}$ be the composition obtained from $\eta$ by changing all $a$’s to $b$, or changing
% all $b$’s to $a$. Then at q = 1, we have
% \LW{Oh!  I see you found a nicer way to write Theorem 4.1
% (than how we had done it in \cite{AyyerMartinWilliams2025})}
% \begin{equation}\label{eq:adjacent}
%     \frac{E^*_\eta(\bfx;1,t)}{E^*_{\widehat{\eta}}(\bfx;1,t)}=\frac{e^*_{\eta'}(\bfx;t)}{e^*_{\widehat{\eta}'}(\bfx;t)}
% \end{equation}
% \end{thm}

\begin{thm}\label{thm:recoloring}
\begin{comment}
Consider two compositions $\mu \in S_n(\lambda)$ and $\eta=\phi(\mu)\in S_n(\kappa)$. Assume that for any $i<j$ such that $\phi(i)=\phi(j)$, all parts of size $i$ in $\mu$ appear to the left of the parts of size $j$ (equivalently, parts that project to the same integer appear in increasing order in $\mu$).
\HBD{For any $\eta$ there is a unique $\mu$ satisfying this condition. Should we extend the notation $\inc$ to compositions and write $\mu=\inc(\eta,\phi).$}
\LW{Just to clarify, this is only unique if we also know $\lambda$, right?  And do you see an advantage in extending the $\inc$ notation?}
\HBD{Yes, you are right! so $\inc(\eta,\phi)$ wouldn't be a good notation. The only reason I thought about this, is that I realized while rereading this section that "the lexicographically smallest" in \cref{def:inc} is the same thing as the condition of \cref{thm:recoloring}. So maybe to make the connection between the two definitions, and instead of extending the notation $\inc$, we can simply reformulate \cref{thm:recoloring} as follows:
\end{comment}
Fix a weakly order-preserving function $\phi:\NN\rightarrow\NN$, and two partitions $\lambda$ and $\kappa$ such that $\phi(\lambda)=\kappa$.
Consider a composition  $\eta\in S_n(\kappa)$. Let $\mu \in S_n(\lambda)$ be the lexicographically smallest composition such that $\phi(\mu)=\eta$ (equivalently, $\mu \in S_n(\lambda)$ such that $\phi(\mu)=\eta$ and parts that project to the same integer appear in increasing order in $\mu$).
 We then have
\begin{equation}\label{eq:adjacent}
    \frac{E^*_\mu(\bfx;1,t)}{E^*_{\eta}(\bfx;1,t)}=\frac{e^*_{\lambda'}(\bfx;t)}{e^*_{\kappa'}(\bfx;t)}.
\end{equation}
\end{thm}


\begin{proof}
\begin{comment}
Assume that $\mu_1 \leq \dots \leq \mu_n$. We then have that $\eta_1 \leq \dots \leq \eta_n$. In this case, we know from \cref{thm:E_partition} that 
$E^*_\mu(\bfx;1,t)=e^*_{\lambda'}(\bfx,t)$ and $E^*_{\eta}(\bfx;1,t)=e^*_{\kappa'}(\bfx,t)$ which gives \cref{eq:adjacent}. We now apply the shape-permuting formula \eqref{shape-permute} to obtain the general results.  \LW{I rewrote the next sentence}
Indeed, to get from an increasing composition $\mu$ to an arbitrary composition
satisfying the conditions of the theorem, we only need to swap two       
parts $\mu_h < \mu_{h+1}$ which project to two different numbers
$\eta_h < \eta_{h+1}$.
%the condition on $\mu$ assures that we only need to invert two parts $\mu_i<\mu_{i+1}$ of $\mu$ such that $\eta_i=\phi(\mu_i)<\phi(\mu_{i+1})=\eta_{i+1}.$ 
Moreover in this case, the shape-permuting operators at $q=1$ have the same behavior on both $E_\mu^*$ and 
 $E^*_\eta$ because $r_j(\mu) = r_j(\eta)$.
\end{comment}

 % \LW{This alternative proof looks good to me.}
 % \HBD{Alternative proof (sketch)}:
We prove the result for any $\lambda$ and $\kappa$, and for any composition $\mu \in S_n(\lambda)$ and $\eta=\phi(\mu)\in S_n(\kappa)$ satisfying the conditions of the theorem. We proceed by induction on the size of $\kappa$.  If $\kappa=(0,\dots,0)$, then $\eta=(0,\dots,0)$, and we have that $\mu$ is weakly increasing ($\mu^{\rev}$ is a partition). As a consequence, applying \cref{thm:E_partition}, we get
\begin{equation}
    \frac{E^*_\mu(\bfx;1,t)}{E^*_{\eta}(\bfx;1,t)}=e^*_{\lambda'}(\bfx;t)=\frac{e^*_{\lambda'}(\bfx;t)}{e^*_{\kappa'}(\bfx;t)}.
\end{equation}
% We can use a similar argument when $\kappa$ is a rectangular partition, i.e. $\kappa=(a,\dots,a)$ for some $a\geq 0$. 

We now fix $\kappa$ and $\lambda$, and we write $\kappa=(a_1^{m_1},\dots,a_\ell^{m_\ell})$, for some $m_i>0$, $a_1>0$ and $a_1>a_2>\dots>a_\ell\geq 0$. We want to prove the result for any $\mu\in S_n(\lambda)$ and $\eta\in S_n(\kappa)$ satisfying the condition of the theorem. It is enough to prove this for $\eta=\kappa$, and the other cases will follow by the shape-permuting formula \cref{prop:shape}. In this case, we know that $\mu$ is of the form 
$\mu=(\mu^{(1)},\dots,\mu^{(\ell)})$, where $\mu^{(i)}$ is a weakly increasing sequence of length $m_i$, and if $i<j$ then all parts of $\mu^{(i)}$ are greater than all parts of $\mu^{(j)}$ (the projection $\phi$ sends then all parts of $\mu^{(i)}$ to $a_i$, for $1\leq i\leq \ell$).
We want to show that the ratio $\frac{E^*_\mu(\bfx;1,t)}{E^*_{\kappa}(\bfx;1,t)}$ can be written in the form of \cref{eq:adjacent}. But using the Knop–Sahi recurrence (\cref{thm:KS_recurrence}), it is enough to prove this result for
$$\frac{E^*_{\nu}(\bfx;1,t)}{E^*_{\zeta}(\bfx;1,t)},$$
where $\nu=(\mu^{(2)},\dots,\mu^{(\ell)},\mu^{(1)}-\mathbf{1})$ and $\zeta=(a_2^{m_2},\dots,a_\ell^{m_\ell},(a_1-1)^{m_1})$.

Now there are two cases:
\begin{itemize}
    \item if $a_1-1=a_2$, we check that in this case we can find a projection $\phi'$ sending $\nu$ on $\zeta$, and that these compositions satisfy the conditions of the theorem. 
    More precisely, we choose $\phi'$ such that it sends all the parts of $\mu^{(i)}$ to $a_i$ for $i\geq 2$, and the parts of $\mu^{(1)}-\mathbf{1}$ to $a_1-1=a_2$ (note that the assumption $a_1-1=a_2$ is crucial for such projection to exist, since the parts of $\mu^{(1)}-\mathbf{1}$ might intersect with the parts of $\mu^{(2)}$). 
    Since $|\zeta|<|\kappa|$,  we can apply the induction hypothesis.

    
    \item if $a_1-1>a_2$, we can write
    $$\frac{E^*_{\nu}(\bfx;1,t)}{E^*_{\zeta}(\bfx;1,t)}=
    \frac{E^*_{\nu}(\bfx;1,t)}{E^*_{(a_2^{m_2},\dots,a_\ell^{m_\ell},a_2^{m_1})}(\bfx;1,t)}\cdotp
    \frac{E_{(a_2^{m_2},\dots,a_\ell^{m_\ell},a_2^{m_1})}(\bfx;1,t)}{E^*_{\zeta}(\bfx;1,t)}$$
    and we apply the induction hypothesis separately on the first term and on the inverse of the second term.\qedhere
\end{itemize}
% We now assume that $|\kappa|>0$.
\end{proof} 


In the following, we fix a weakly order-preserving function $\phi:\NN\rightarrow\NN$, and two partitions $\lambda$ and $\kappa$ such that $\phi(\lambda)=\kappa$.


\begin{definition}\label{def:inc}
Fix $\phi$, $\lambda$ and $\kappa$ as above.
Let $\inc(\lambda,\phi)$ be the lexicographically smallest $\rho\in S_n(\lambda)$ such that $\phi(\rho)=\kappa$. 
\end{definition}

\begin{thm}\label{thm:E_inc_E}
    Fix $\kappa$ and $\lambda$ as in \cref{def:inc}. At $q=1$, $E^*_{\inc(\lambda,\phi)}$
    % \LW{$E^*_{\inc(\lambda, \phi)}$?}
    is a symmetric function multiple of $E^*_\kappa=F^*_\kappa$.
    \end{thm}
%    \LW{Do we actually need \cref{thm:standardization} here?  I was thinking of the example 
%$\phi:(6,6,4,3,3,1) \to %(4,4,4,1,1,1)$ so that 
%$\rho=\inc(\lambda,\phi) = %(4,4,6,1,3,3)$
%and it seems that from 
%$$\frac{E^*_{446133}}%{E^*_{444111}} = 
%\frac{E^*_{446133}}{E^*_{444133}} \cdot\frac{E^*_{444133}}{E^*_{444111}}$$
%we can just use \cref{thm:adjacent} to get what we want.}
%\HBD{I think that with the old version of \cref{thm:recoloring} (the one with adjacent parts),  we need some form of standardization here, with examples like $\phi:(4,3,2,1)\rightarrow (2,2,1,1)$, right?}
 
\begin{proof}
    We apply \cref{thm:recoloring} with $\eta=\kappa$ and $\mu=\inc(\lambda,\phi)$.
\end{proof}


Recall that for $\eta\in S_n(\kappa)$, 
\begin{equation}\label{eq:G}
G^*_\eta(\bfx;t):=\sum_{\rho\in S_n(\lambda):\, \phi(\rho)=\eta} F^*_\rho(\bfx;1,t),
\end{equation}
and that $G_{\eta}$ is the top homogeneous part of $G^*_\eta$.

\begin{prop}[{\cite[Proposition 4.16]{AyyerMartinWilliams2025}}]\label{prop:hom_G}
    We have 
$G_\kappa(\bfx;t)=E_{\inc(\lambda,\phi)}(\bfx;1,t).$
\end{prop}
We then deduce the following theorem for interpolation polynomials.
\begin{prop}\label{prop:G_E}
    We have 
$G^*_\kappa(\bfx;t)=E^*_{\inc(\lambda,\phi)}(\bfx;1,t).$
\end{prop}
\begin{proof} By definition $G^*_\kappa$ is a linear combination of $F^*_\nu(\bfx;1,t)$ for $\nu\in S_n(\lambda)$. But we know from \cref{prop:Vstar} that for $\nu\in S_n(\lambda)$, $F^*_\nu(\bfx;1,t)$ is a linear combination of $E^*_\rho(\bfx;1,t)$ for $\rho\in S_n(\lambda)$. 
% \LW{Better to replace these last two instances of $\kappa$ by another Greek letter, since our proposition, $\kappa$ is fixed}. 
We find the coefficients of this expansion by taking the top homogeneous part and using \cref{prop:hom_G}.
\end{proof}


% The following is an analogue of \cite[Theorem 4.17]{AyyerMartinWilliams2025}. It is a consequence of \cref{thm:adjacent} and \cref{thm:standardization}.
\begin{proof}[Proof of \cref{thm:weak-reordering}]
We want to prove that for all $\eta\in S_n(\kappa)$, we have at $q=1$
    $$\frac{G^*_\eta(\bfx;t)}{P^*_\lambda(\bfx;1,t)}=\frac{F^*_\eta(\bfx;1,t)}{P^*_\kappa(\bfx;1,t)}.$$
    We know from \cref{prop:G_E} that $G_\kappa^*(\bfx;t)=E^*_{\inc(\lambda,\phi)}(\bfx;1,t)$. Using \cref{thm:E_inc_E}, we get that 
    \begin{equation}\label{eq:G-F1}
      G_\kappa^*(\bfx;t)
      =h\cdotp E_\kappa^*(\bfx;1,t)
      =h\cdotp F_\kappa^*(\bfx;1,t),  
    \end{equation}
    for some symmetric function $h$.
    Now by \cref{prop:T_i-F_star}, for any
    $\eta \in S_n(\kappa)$, we have 
    %    we can obtain any $F_\eta^*$ from $F_\kappa^*$ by applying a Hecke operators $T_\sigma$ (see ). More precisely 
    $F^*_\eta=T_\sigma \cdotp F_\kappa^*,$
    where $\sigma$ is the shortest permutation sending $\kappa$ to $\eta$. One then checks that we also have $G^*_\eta=T_\sigma \cdotp G_\kappa^*$ (see the proof of \cite[Lemma 4.12]{AyyerMartinWilliams2025} for more details).
    Now applying $T_\sigma$ to \cref{eq:G-F1} we get that 
    $ G_\eta^*(\bfx;t)=h\cdotp F_\eta^*(\bfx;1,t).$
%    for any $\eta\in S_n(\kappa).$
    Summing this equation over all $\eta$, we get
    $$P_\lambda^*(\bfx;1,t)=h\cdotp P_\kappa^*(\bfx;1,t),$$
    where we have used \cref{eq:G} and \cref{prop:symmetrization}. We then get the desired claim
    \begin{equation*}
    \frac{P_\lambda^*(\bfx;1,t)}{P_\kappa^*(\bfx;1,t)}=h=\frac{G_\eta^*(\bfx;t)}{F_\eta^*(\bfx;1,t)}.\qedhere    
    \end{equation*}
%    for any $\eta\in S_n(\kappa)$ as claimed.
\end{proof}


\section{\texorpdfstring{$t$}{t}-interpolation Schur polynomials}\label{app:t_int_Schur}
% \LW{appendix?}
Throughout this section we fix a positive integer $n$,
and write 
$g(\bfx)=g(x_1,\dots,x_n)$.
Recall that the $t$-interpolation Schur polynomial $s_\lambda^*$ is the symmetric polynomial defined by
    $$s^*_\lambda(\bfx;t)=P_\lambda(\bfx;t,t).$$


We let $\SSYT(\lambda;n)$ denote the set of \emph{semistandard Young tableaux} of shape $\lambda$ filled with values in $\llbracket n\rrbracket$. For a cell $\Box=(i,j)$ in the Young diagram of $\lambda$ in  row $i$ and column $j$, we let $c(\Box)=j-i$ denote the \emph{content} of $\Box$.

With the specialization $q=t$, Okounkov's formula for interpolation Macdonald polynomials \cite[Equation (1.4)]{Okounkov1998} can be written as follows.

\begin{thm}[\cite{Okounkov1998}]\label{thm:Okounkov}
For any partition $\lambda=(\lambda_1,\dots,\lambda_n)$, we have
$$s^*_\lambda(\bfx;t)=\sum_{T\in \SSYT(\lambda)}\prod_{\Box\in \lambda}(x_{T(\Box)}-t^{c(\Box)+T(\Box)-n}).$$
\end{thm}

For $n\geq k\geq 0$ and $\ell\in \ZZ$, we define
$$e_{k,\ell}^*(\bfx;t):=\sum_{\begin{subarray}{c}S\subseteq \llbracket n \rrbracket\\  \#S=k\end{subarray}}\prod_{i\in S}\left(x_i-\frac{t^{\#S^c\cap \llbracket i-1\rrbracket}}{t^{\ell}}\right),$$
so that $e^*_k(\bfx;t)=e^*_{k,n-1}(\bfx;t)$
and also
\begin{equation}\label{eq:identity}
e^*_{k,n-2}(tx_1,\dots,tx_n;t) = t^k e^*_{k,n-1}(x_1,\dots,x_n;t).    
\end{equation}


We then have the following $t$-deformation of the dual Jacob–Trudi identity. It follows from \cref{thm:Okounkov}, together with the results of \cite{Macdonald1992} on the 6th variation of Schur functions (more precisely, with the notation \cite[Section6]{Macdonald1992}, we set $a_i=-t^{i-n}$).

\begin{prop}[\cite{Macdonald1992}]
    Let $\lambda=(\lambda_1,\dots,\lambda_{n})$ be a partition. We have
    \begin{equation}\label{eq:int_dual_JT}
      s^*_\lambda(\bfx;t)=\det\left(e^*_{\lambda'_i-i+j, n-j}(\bfx;t)\right)_{1\leq i,j\leq n}.  
    \end{equation}
\end{prop}

\begin{remark}
    Note that the top homogeneous part of \cref{eq:int_dual_JT} gives the dual Jacobi–Trudi identity (see \textit{e.g} \cite[Corollary 7.16.2]{Stanley:EC2}). As in the homogeneous case, this result can also be obtained using the Lindström–Gessel–Viennot lemma.
\end{remark}

If $\lambda$
%=(2^a,1^b)$ 
is a 2-column partition, then \cref{eq:identity} and \cref{eq:int_dual_JT} give the following.
\begin{lem}\label{lem:2column}
Let $\lambda=(2^a,1^b)$.  Then we have 
\begin{align*}
s^*_{\lambda}(t\bfx;t)&=e^*_{a+b,n-1}(t\bfx;t)e^*_{a,n-2}(t\bfx;t)-e^*_{a+b+1,n-2}(t\bfx;t)e^*_{a-1,n-1}(t\bfx;t)\\
&=t^a e^*_{a+b}(t\bfx;t) e^*_a(\bfx;t) - t^{a+b+1} e^*_{a+b+1}(\bfx;t) e^*_{a-1}(t\bfx;t).
\end{align*}
% where $t\bfx=(tx_1,\dots,tx_n)$.
% \HBD{Should the left-hand side be evaluated at $t\bfx$ instead of $\bfx$:
% \begin{align*}
% s^*_{\lambda}(t\bfx;t)&=e^*_{a+b,n-1}(t\bfx;t)e^*_{a,n-2}(t\bfx;t)-e^*_{a+b+1,n-2}(t\bfx;t)e^*_{a-1,n-1}(t\bfx;t)\\
% &=t^a e^*_{a+b}(tx_1,\dots,tx_n) e^*_a(x_1,\dots,x_n) - t^{a+b+1} e^*_{a+b+1}(x_1,\dots,x_n) e^*_{a-1}(tx_1,\dots,tx_n).
% \end{align*}
% }
% \LW{Yes!}
\end{lem}
\end{appendices}


\fi
%%%%%%%%%%%%%%%%%%%%%%%%%%%%%%%%%%%%%%%%%
%%%%%%%%%%%%%%%%%%%%%%%%%%%%%%%

\bibliographystyle{amsalpha}
\bibliography{biblio.bib}
   
\end{document}
